\chapter{Derivatives}
\label{chap:pet-derivatives}

Faithful transcription of Petersen, \emph{Riemannian Geometry} (GTM 171, 3rd ed.),
Chapter 2. The chapter develops the various notions of derivative of a tensor on
a manifold: the Lie derivative (defined via flows, and characterized
algebraically through the Lie bracket and a Leibniz rule), the affine and
Riemannian connections and covariant differentiation (culminating in the
Fundamental Theorem of Riemannian Geometry, which characterizes the Levi-Civita
connection by torsion-freeness and compatibility with the metric), the unified
algebraic treatment of Lie and covariant derivatives as "natural derivations"
coming from the action of $\mathrm{GL}(V)$ and $\mathrm{End}(V)$ on the tensor
algebra, and the classical coordinate (Christoffel symbol) formulas for the
connection. Along the way the chapter defines, and relates to each other, the
gradient, Hessian, divergence and Laplacian of a function, and closes with an
extensive list of exercises.

\section{Lie Derivatives}

\subsection{Directional Derivatives}

\begin{definition}[directional derivative]
  \label{def:pet-ch2-directional-derivative}
  \lean{PetersenLib.directionalDerivative,PetersenLib.gradient}
  \leanok
  Given a function $f:M\to\mathbb{R}$ and a vector field $Y$ on $M$, the
  \emph{directional derivative} of $f$ in the direction of $Y$ is written
  interchangeably as
  \[
    \nabla_Yf=D_Yf=L_Yf=df(Y)=Y(f).
  \]
  If $M$ carries a Riemannian metric $g$, the \emph{gradient} $\operatorname{grad}f=\nabla f$
  is the vector field with $g(v,\nabla f)=df(v)$ for all $v\in TM$; in local
  coordinates $\nabla f=g^{ij}\partial_i(f)\partial_j$, where $[g^{ij}]$ is the inverse of
  $[g_{ij}]$.
  \source{petersen:page-0060}
\end{definition}

\begin{remark}[the gradient is not metric-independent]
  \label{rem:pet-ch2-gradient-not-invariant}
  \lean{PetersenLib.gradient_not_coordinate_invariant}
  \uses{def:pet-ch2-directional-derivative}
  On $\mathbb{R}^n$ with Cartesian coordinates, $\nabla f=g^{ij}\partial_i(f)\partial_j=\sum_i\partial_i(f)\partial_i$,
  but this expression is not invariantly defined: in polar coordinates on
  $\mathbb{R}^2$, $\partial_r(f)\partial_r+\partial_\theta(f)\partial_\theta$ (formed the same naive way) does
  \emph{not} equal the true gradient. A rule of thumb is that invariantly
  defined objects satisfy the Einstein summation convention with all indices
  paired: $df=\partial_i(f)dx^i$ is invariant, $g=g_{ij}dx^idx^j$ and
  $\nabla f=g^{ij}\partial_i(f)\partial_j$ are invariant (given a metric), but a bare
  $\partial_i(f)\partial_i$ is not.
  \source{petersen:page-0060}
\end{remark}

\subsection{Lie Derivatives}

\begin{definition}[Lie derivative of a function]
  \label{def:pet-ch2-lie-derivative-function}
  \lean{PetersenLib.lieDerivativeFunction}
  \leanok
  Let $X$ be a vector field and $F^t$ its locally defined flow on $M$, so
  $\gamma(t)=F^t(p)$ is the integral curve of $X$ through $p$ at $t=0$. For
  $f:M\to\mathbb{R}$, the \emph{Lie derivative} $L_Xf$ is the vector field's first-order
  effect on $f$ under the flow:
  \[
    f\big(F^t(p)\big)=f(p)+t(L_Xf)(p)+o(t),
    \qquad
    (L_Xf)(p)=\lim_{t\to0}\frac{f(F^t(p))-f(p)}{t},
  \]
  or without reference to $p$, $f\circ F^t=f+tL_Xf+o(t)$. Thus $L_Xf=D_Xf=df(X)$
  (\cref{def:pet-ch2-directional-derivative}).
  \source{petersen:page-0061}
\end{definition}

\begin{definition}[Lie derivative of a vector field]
  \label{def:pet-ch2-lie-derivative-vector-field}
  \lean{PetersenLib.lieDerivativeVectorField}
  \leanok
  \uses{def:pet-ch2-lie-derivative-function}
  For a vector field $Y$, the value $Y|_{F^t(p)}$ cannot be compared directly to
  $Y|_p$; instead consider the curve $t\mapsto DF^{-t}\big(Y|_{F^t(p)}\big)\in T_pM$.
  Its first-order Taylor expansion at $t=0$,
  \[
    DF^{-t}\big(Y|_{F^t(p)}\big)=Y|_p+t(L_XY)|_p+o(t),
  \]
  defines the \emph{Lie derivative} $(L_XY)|_p\in T_pM$, equivalently
  $(L_XY)|_p=\lim_{t\to0}\frac{DF^{-t}(Y|_{F^t(p)})-Y|_p}{t}$.
  \source{petersen:page-0061}
\end{definition}

\begin{proposition}[Prop. 2.1.1 — Lie derivative of a vector field is the bracket]
  \label{prop:pet-ch2-lie-derivative-bracket}
  \lean{PetersenLib.lieDerivative_vectorField_eq_bracket}
  \leanok
  \uses{def:pet-ch2-lie-derivative-vector-field}
  If $X,Y$ are vector fields on $M$, then $L_XY=[X,Y]$.
  \source{petersen:page-0061,petersen:page-0062}
\end{proposition}
\begin{proof}
  \uses{prop:pet-ch2-lie-derivative-bracket}
  \leanok
  Writing the defining expansion as $Y|_{F^t}-DF^t(Y)=tDF^t(L_XY)+o(t)$, take the
  directional derivative of a function $f$ along $Y|_{F^t}-DF^t(Y)$:
  \[
    D_{Y|_{F^t}-DF^t(Y)}f
    = D_{Y|_{F^t}}f-D_{DF^t(Y)}f
    = (D_Yf)\circ F^t-D_Y\big(f\circ F^t\big).
  \]
  Expanding both terms to first order in $t$,
  \[
    (D_Yf)\circ F^t = D_Yf+tD_XD_Yf+o(t),
    \qquad
    D_Y(f\circ F^t) = D_Y(f+tD_Xf+o(t)) = D_Yf+tD_YD_Xf+o(t),
  \]
  so
  \[
    D_{Y|_{F^t}-DF^t(Y)}f = t\big(D_XD_Yf-D_YD_Xf\big)+o(t) = tD_{[X,Y]}f+o(t).
  \]
  On the other hand $D_{Y|_{F^t}-DF^t(Y)}f = tD_{DF^t(L_XY)}f+o(t)$, and as $t\to0$,
  $DF^t(L_XY)\to L_XY$, so comparing first-order terms as directional
  derivatives of an arbitrary $f$ gives $L_XY=[X,Y]$.
\end{proof}

\begin{definition}[Lie derivative of a $(0,k)$-tensor]
  \label{def:pet-ch2-lie-derivative-tensor}
  \lean{PetersenLib.lieDerivativeTensor}
  \leanok
  \uses{def:pet-ch2-lie-derivative-vector-field}
  For a $(0,k)$-tensor $T$, the pullback $(F^t)^*T$ satisfies
  \[
    \big((F^t)^*T\big)(Y_1,\dots,Y_k)=T\big(DF^t(Y_1),\dots,DF^t(Y_k)\big)
    =T(Y_1,\dots,Y_k)+t(L_XT)(Y_1,\dots,Y_k)+o(t),
  \]
  defining the \emph{Lie derivative} $L_XT$, equivalently
  $(L_XT)(Y_1,\dots,Y_k)=\lim_{t\to0}\frac{(F^t)^*T-T}{t}(Y_1,\dots,Y_k)$.
  \source{petersen:page-0062}
\end{definition}

\begin{proposition}[Prop. 2.1.2 — algebraic formula for $L_XT$]
  \label{prop:pet-ch2-lie-derivative-tensor-formula}
  \lean{PetersenLib.lieDerivativeTensor_formula}
  \leanok
  \uses{def:pet-ch2-lie-derivative-tensor,prop:pet-ch2-lie-derivative-bracket}
  If $X$ is a vector field and $T$ a $(0,k)$-tensor, then
  \[
    (L_XT)(Y_1,\dots,Y_k)
    = D_X\big(T(Y_1,\dots,Y_k)\big)-\sum_{i=1}^kT(Y_1,\dots,L_XY_i,\dots,Y_k).
  \]
  \source{petersen:page-0062,petersen:page-0063}
\end{proposition}
\begin{proof}
  \uses{prop:pet-ch2-lie-derivative-tensor-formula}
  \leanok
  It suffices to treat $k=1$; the general case follows the same pattern with
  more indices. From \cref{def:pet-ch2-lie-derivative-vector-field},
  $Y|_{F^t}=DF^t(Y)+tDF^t(L_XY)+o(t)$, so
  \[
    \big((F^t)^*T\big)(Y) = T(DF^t(Y))
    = T\big(Y|_{F^t}-tDF^t(L_XY)\big)+o(t)
    = T(Y)\circ F^t - tT(DF^t(L_XY))+o(t)
  \]
  \[
    = T(Y)+tD_X\big(T(Y)\big)-tT\big(DF^t(L_XY)\big)+o(t).
  \]
  Hence
  \[
    (L_XT)(Y)=\lim_{t\to0}\frac{\big((F^t)^*T\big)(Y)-T(Y)}{t}
    = \lim_{t\to0}\Big(D_X(T(Y))-T\big(DF^t(L_XY)\big)\Big)
    = D_X(T(Y))-T(L_XY),
  \]
  using $DF^t(L_XY)\to L_XY$ as $t\to0$.
\end{proof}


\begin{definition}[product of tensors]
  \label{def:pet-ch2-tensor-product}
  \lean{PetersenLib.TensorOperator.mul}
  \leanok
  \uses{def:pet-ch2-lie-derivative-tensor}
  For a $(0,k_1)$-tensor $T_1$ and a $(0,k_2)$-tensor $T_2$, their
  \emph{product} is the $(0,k_1+k_2)$-tensor
  \[
    (T_1\cdot T_2)(X_1,\dots,X_{k_1},Y_1,\dots,Y_{k_2})
    =T_1(X_1,\dots,X_{k_1})\cdot T_2(Y_1,\dots,Y_{k_2}).
  \]
  \source{petersen:page-0063}
\end{definition}

\begin{proposition}[Prop. 2.1.3 — product rule]
  \label{prop:pet-ch2-lie-derivative-product-rule}
  \lean{PetersenLib.lieDerivative_product_rule}
  \leanok
  \uses{prop:pet-ch2-lie-derivative-tensor-formula}
  If $T_1,T_2$ are $(0,k_i)$-tensors, then
  $L_X(T_1\cdot T_2)=(L_XT_1)\cdot T_2+T_1\cdot(L_XT_2)$, where
  $(T_1\cdot T_2)(X_1,\dots,X_{k_1},Y_1,\dots,Y_{k_2})
  =T_1(X_1,\dots,X_{k_1})\cdot T_2(Y_1,\dots,Y_{k_2})$.
  \source{petersen:page-0063}
\end{proposition}
\begin{proof}
  \uses{prop:pet-ch2-lie-derivative-product-rule}
  \leanok
  By \cref{prop:pet-ch2-lie-derivative-tensor-formula},
  $L_X(T_1\cdot T_2)$ evaluated on $(X_1,\dots,X_{k_1},Y_1,\dots,Y_{k_2})$ is the
  directional derivative $D_X$ of $T_1(X_\bullet)\cdot T_2(Y_\bullet)$ minus the
  sum, over each slot, of the tensor with $L_X$ applied to that slot's vector
  field. Applying the ordinary product rule for $D_X$ to
  $T_1(X_\bullet)\cdot T_2(Y_\bullet)$ and regrouping the sum over the
  $X$-slots against $T_2$ and the sum over the $Y$-slots against $T_1$
  reproduces exactly $(L_XT_1)\cdot T_2+T_1\cdot(L_XT_2)$, again by
  \cref{prop:pet-ch2-lie-derivative-tensor-formula} applied to each factor.
\end{proof}

\begin{proposition}[Prop. 2.1.4 — Lie derivative in the direction $fX$]
  \label{prop:pet-ch2-lie-derivative-scaling}
  \lean{PetersenLib.lieDerivative_scale_direction}
  \leanok
  \uses{prop:pet-ch2-lie-derivative-tensor-formula}
  If $T$ is a $(0,k)$-tensor and $f:M\to\mathbb{R}$, then
  \[
    (L_{fX}T)(Y_1,\dots,Y_k)
    = f(L_XT)(Y_1,\dots,Y_k)+\sum_{i=1}^k(L_{Y_i}f)\,T(Y_1,\dots,X,\dots,Y_k),
  \]
  with $X$ replacing $Y_i$ in the $i$th slot.
  \source{petersen:page-0063,petersen:page-0064}
\end{proposition}
\begin{proof}
  \uses{prop:pet-ch2-lie-derivative-scaling}
  \leanok
  By \cref{prop:pet-ch2-lie-derivative-tensor-formula},
  \[
    (L_{fX}T)(Y_1,\dots,Y_k)
    = D_{fX}\big(T(Y_1,\dots,Y_k)\big)-\sum_{i=1}^kT(Y_1,\dots,[fX,Y_i],\dots,Y_k)
  \]
  \[
    = fD_X\big(T(Y_1,\dots,Y_k)\big)-\sum_{i=1}^kT(Y_1,\dots,f[X,Y_i]-(L_{Y_i}f)X,\dots,Y_k)
  \]
  using $[fX,Y_i]=f[X,Y_i]-(D_{Y_i}f)X$. Distributing the tensor over the
  sum inside each slot and using multilinearity of $T$ splits this into
  \[
    f\Big(D_X\big(T(Y_1,\dots,Y_k)\big)-\sum_iT(Y_1,\dots,[X,Y_i],\dots,Y_k)\Big)
    +\sum_i(L_{Y_i}f)\,T(Y_1,\dots,X,\dots,Y_k),
  \]
  and the first bracket is $f(L_XT)(Y_1,\dots,Y_k)$ again by
  \cref{prop:pet-ch2-lie-derivative-tensor-formula}.
\end{proof}

\begin{lemma}[Lem. 2.1.5 — locality at a zero of $X$]
  \label{lem:pet-ch2-lie-derivative-local-vanishing}
  \lean{PetersenLib.lieDerivative_local_at_zero}
  \leanok
  \uses{prop:pet-ch2-lie-derivative-tensor-formula}
  If a vector field $X$ vanishes at $p$, then $(L_XT)|_p$ depends only on the
  value $T|_p$ of the $(0,k)$-tensor $T$ at $p$.
  \source{petersen:page-0064}
\end{lemma}
\begin{proof}
  \uses{lem:pet-ch2-lie-derivative-local-vanishing}
  \leanok
  By \cref{prop:pet-ch2-lie-derivative-tensor-formula},
  $(L_XT)(Y_1,\dots,Y_k)=D_X(T(Y_1,\dots,Y_k))-\sum_iT(Y_1,\dots,L_XY_i,\dots,Y_k)$.
  At $p$ with $X|_p=0$ the directional derivative term $D_X(T(Y_\bullet))|_p$
  vanishes (it only involves $X|_p$), leaving
  $(L_XT)(Y_1,\dots,Y_k)|_p=-\sum_iT(Y_1,\dots,L_XY_i,\dots,Y_k)|_p$, an
  expression built from $T|_p$ (evaluated on the vectors $Y_j|_p$ and
  $(L_XY_i)|_p$) alone.
\end{proof}

\begin{definition}[iterated Lie derivative $(L_XL)_YT$]
  \label{def:pet-ch2-generalized-lie-derivative}
  \lean{PetersenLib.generalizedLieDerivative}
  \leanok
  \uses{def:pet-ch2-lie-derivative-tensor}
  For vector fields $X,Y$ and a tensor (or multilinear map on vector fields)
  $T$,
  \[
    (L_XL)_YT := L_X(L_YT)-L_{L_XY}T-L_Y(L_XT) = [L_X,L_Y]T-L_{[X,Y]}T.
  \]
  \source{petersen:page-0064}
\end{definition}

\begin{proposition}[Prop. 2.1.6 — the Generalized Jacobi Identity]
  \label{prop:pet-ch2-generalized-jacobi-identity}
  \lean{PetersenLib.generalizedJacobiIdentity}
  \leanok
  \uses{def:pet-ch2-generalized-lie-derivative}
  For all vector fields $X,Y$ and tensors $T$, $(L_XL)_YT=0$.
  \source{petersen:page-0065,petersen:page-0066}
\end{proposition}
\begin{proof}
  \uses{prop:pet-ch2-generalized-jacobi-identity}
  \leanok
  \emph{Functions.} $(L_XL)_Yf=[L_X,L_Y]f-L_{[X,Y]}f=[D_X,D_Y]f-D_{[X,Y]}f=0$,
  the last equality being the definition of the Lie bracket as a commutator of
  directional derivatives.

  \emph{Vector fields.} For $T=Z$ a vector field,
  using \cref{prop:pet-ch2-lie-derivative-bracket},
  \[
    (L_XL)_YZ=[L_X,L_Y]Z-L_{[X,Y]}Z
    =[X,[Y,Z]]-[Y,[X,Z]]-[[X,Y],Z]
    =[X,[Y,Z]]+[Z,[X,Y]]+[Y,[Z,X]]=0
  \]
  by the ordinary Jacobi identity for the Lie bracket.

  \emph{One-forms.} For $T=\omega$ a one-form, first note that
  \[
    \big([L_X,L_Y]\omega\big)(Z)=[L_X,L_Y]\big(\omega(Z)\big)-\omega\big([L_X,L_Y]Z\big).
  \]
  Indeed, by \cref{prop:pet-ch2-lie-derivative-tensor-formula},
  \[
    \big(L_X(L_Y\omega)\big)(Z)
    = L_X\big((L_Y\omega)(Z)\big)-(L_Y\omega)(L_XZ)
    = L_X\big(L_Y(\omega(Z))\big)-L_X\big(\omega(L_YZ)\big)
    -L_Y\big(\omega(L_XZ)\big)+\omega(L_YL_XZ),
  \]
  and symmetrically
  \[
    \big(L_Y(L_X\omega)\big)(Z)
    = L_Y\big(L_X(\omega(Z))\big)-L_Y\big(\omega(L_XZ)\big)-L_X\big(\omega(L_YZ)\big)+\omega(L_XL_YZ).
  \]
  Subtracting, all cross terms cancel and
  \[
    \big([L_X,L_Y]\omega\big)(Z)=(L_X(L_Y\omega))(Z)-(L_Y(L_X\omega))(Z)
    = [L_X,L_Y]\big(\omega(Z)\big)-\omega\big([L_X,L_Y]Z\big).
  \]
  Consequently
  \[
    \big((L_XL)_Y\omega\big)(Z)
    = \big([L_X,L_Y]\omega\big)(Z)-\big(L_{[X,Y]}\omega\big)(Z)
    = [L_X,L_Y]\big(\omega(Z)\big)-\omega\big([L_X,L_Y]Z\big)
    -L_{[X,Y]}\big(\omega(Z)\big)+\omega\big(L_{[X,Y]}Z\big).
  \]
  The function terms $[L_X,L_Y](\omega(Z))-L_{[X,Y]}(\omega(Z))$ vanish by the
  function case, and $-\omega([L_X,L_Y]Z)+\omega(L_{[X,Y]}Z)
  =-\omega\big((L_XL)_YZ\big)$ vanishes by the vector-field case; hence
  $(L_XL)_Y\omega=0$.

  \emph{General tensors.} A general $(0,k)$-tensor is (locally, on a spanning
  set) a sum of tensor products of functions, vector fields, and one-forms.
  Since $L_X$ acts as a derivation on tensor products
  (\cref{prop:pet-ch2-lie-derivative-product-rule}), $(L_XL)_Y$ inherits a
  Leibniz-type expansion over such products, and each factor contributes zero
  by the three cases above; hence $(L_XL)_YT=0$ for all tensors $T$.
\end{proof}

\begin{definition}[exterior derivative via Lie derivatives]
  \label{def:pet-ch2-exterior-derivative-lie-formula}
  \lean{PetersenLib.exteriorDerivative_lieFormula}
  \leanok
  \uses{def:pet-ch2-lie-derivative-tensor}
  For a $k$-form $\omega$,
  \[
    d\omega(X_0,\dots,X_k)
    =\frac12\sum_{i=0}^k(-1)^i(L_{X_i}\omega)(X_0,\dots,\hat X_i,\dots,X_k)
    +\frac12\sum_{i=0}^k(-1)^iL_{X_i}\big(\omega(X_0,\dots,\hat X_i,\dots,X_k)\big).
  \]
  For a $1$-form this specializes to
  $d\omega(X,Y)=D_X(\omega(Y))-D_Y(\omega(X))-\omega([X,Y])$.
  \source{petersen:page-0066}
\end{definition}

\subsection{Lie Derivatives and the Metric}

\begin{definition}[Hessian via the Lie derivative]
  \label{def:pet-ch2-hessian-lie-derivative}
  \lean{PetersenLib.hessianLieDerivative}
  \leanok
  \uses{def:pet-ch2-lie-derivative-tensor,def:pet-ch2-directional-derivative}
  The \emph{Hessian} of $f:(M,g)\to\mathbb{R}$ is the $(0,2)$-tensor
  $\operatorname{Hess}f(X,Y)=\frac12(L_{\nabla f}g)(X,Y)$.
  \source{petersen:page-0066}
\end{definition}

\begin{proposition}[Hessian at a critical point is metric-independent]
  \label{prop:pet-ch2-hessian-critical-point}
  \lean{PetersenLib.hessian_criticalPoint_metricIndependent}
  \uses{def:pet-ch2-hessian-lie-derivative}
  If $df|_p=0$, choose coordinates $x^i$ around $p$ with $g_{ij}|_p=\delta_{ij}$.
  Then $\operatorname{Hess}f|_p$ computed from $g$ equals $\operatorname{Hess}f|_p$
  computed from the flat metric $\delta_{ij}dx^idx^j$ in the same coordinates.
  \source{petersen:page-0066,petersen:page-0067}
\end{proposition}
\begin{proof}
  \uses{prop:pet-ch2-hessian-critical-point}
  Since $df|_p=0$, also $\nabla f|_p=0$, so
  \[
    L_{\nabla f}\big(g_{ij}dx^idx^j\big)\big|_p
    = L_{\nabla f}(g_{ij})|_p\cdot dx^idx^j+g_{ij}|_p\,L_{\nabla f}(dx^i)dx^j+g_{ij}|_p\,dx^iL_{\nabla f}(dx^j).
  \]
  By \cref{lem:pet-ch2-lie-derivative-local-vanishing}, applied to the field
  $\nabla f$ vanishing at $p$, the coefficient function $g_{ij}$ (a $(0,0)$
  contribution) does not affect the value at $p$ beyond its value
  $g_{ij}|_p=\delta_{ij}$, and using this substitution,
  \[
    L_{\nabla f}(g_{ij}dx^idx^j)|_p
    = \delta_{ij}L_{\nabla f}(dx^i)dx^j+\delta_{ij}dx^iL_{\nabla f}(dx^j)
    = L_{\nabla f}(\delta_{ij}dx^idx^j)|_p,
  \]
  i.e.\ $L_{\nabla f}g|_p=L_{\nabla f}(\delta_{ij}dx^idx^j)|_p$, so
  $\operatorname{Hess}f|_p=\frac12(L_{\nabla f}g)|_p$ agrees with the Hessian
  computed from the Euclidean metric $\delta_{ij}dx^idx^j$ in the same fixed
  coordinate system.
\end{proof}

\begin{definition}[divergence via the Lie derivative]
  \label{def:pet-ch2-divergence-lie-derivative}
  \lean{PetersenLib.divergenceLieDerivative}
  \leanok
  \uses{def:pet-ch2-lie-derivative-tensor}
  The \emph{divergence} of a vector field $X$ is the function $\operatorname{div}X$
  measuring how the volume form changes along the flow of $X$:
  $L_X\,\mathrm{vol}=(\operatorname{div}X)\,\mathrm{vol}$.
  \source{petersen:page-0067}
\end{definition}

\begin{remark}[$L_X\mathrm{vol}$ is exact]
  \label{rem:pet-ch2-lie-derivative-volume-exact}
  \lean{PetersenLib.lieDerivative_volume_exact}
  \uses{def:pet-ch2-divergence-lie-derivative}
  $L_X\mathrm{vol}=d\,i_X\mathrm{vol}$, where $i_XT$ evaluates $T$ on $X$ in
  its first variable.
  \source{petersen:page-0067}
\end{remark}

\begin{definition}[Laplacian]
  \label{def:pet-ch2-laplacian}
  \lean{PetersenLib.laplacian}
  \leanok
  \uses{def:pet-ch2-divergence-lie-derivative,def:pet-ch2-directional-derivative}
  The \emph{Laplacian} of $f$ is $\Delta f=\operatorname{div}\nabla f$.
  \source{petersen:page-0067}
\end{definition}

\begin{proposition}[divergence as a trace of $L_Xg$]
  \label{prop:pet-ch2-divergence-trace-formula}
  \lean{PetersenLib.divergence_trace_formula}
  \leanok
  \uses{def:pet-ch2-divergence-lie-derivative}
  For a positively oriented orthonormal frame $E_1,\dots,E_n$,
  $\operatorname{div}X=\sum_i\frac12(L_Xg)(E_i,E_i)$; consequently
  $\Delta f=\operatorname{div}\nabla f=\sum_i\operatorname{Hess}f(E_i,E_i)$.
  \source{petersen:page-0067}
\end{proposition}
\begin{proof}
  \uses{prop:pet-ch2-divergence-trace-formula}
  \leanok
  Using \cref{prop:pet-ch2-lie-derivative-tensor-formula} on
  $\mathrm{vol}(E_1,\dots,E_n)$,
  \[
    \operatorname{div}X = (L_X\mathrm{vol})(E_1,\dots,E_n)
    = L_X\big(\mathrm{vol}(E_1,\dots,E_n)\big)-\sum_i\mathrm{vol}(E_1,\dots,L_XE_i,\dots,E_n)
    = -\sum_ig(L_XE_i,E_i),
  \]
  since $\mathrm{vol}(E_1,\dots,E_n)\equiv1$ is constant and, in the sum,
  replacing one frame slot $E_i$ by the general vector $L_XE_i$ evaluates the
  volume form as the component of $L_XE_i$ along that direction, namely
  $g(L_XE_i,E_i)$. Finally
  \[
    -\sum_ig(L_XE_i,E_i)
    = \frac12\sum_i\Big(D_X\big(g(E_i,E_i)\big)-g(L_XE_i,E_i)-g(E_i,L_XE_i)\Big)
    = \sum_i\frac12(L_Xg)(E_i,E_i),
  \]
  using $D_X(g(E_i,E_i))=0$ (since $g(E_i,E_i)\equiv1$) and
  \cref{prop:pet-ch2-lie-derivative-tensor-formula}.
\end{proof}

\begin{proposition}[Hessian agrees with the Euclidean Hessian]
  \label{prop:pet-ch2-hessian-euclidean-agreement}
  \lean{PetersenLib.hessian_euclidean_agreement}
  \uses{def:pet-ch2-hessian-lie-derivative}
  For $f:\mathbb{R}^n\to\mathbb{R}$ with the canonical metric, the Hessian of
  \cref{def:pet-ch2-hessian-lie-derivative} equals the classical Hessian:
  $L_{\nabla f}(\delta_{ij}dx^idx^j)=2\sum_i\partial_{ij}f\,dx^idx^j$, i.e.
  $\operatorname{Hess}f=\sum_i\partial_{ij}f\,dx^idx^j$.
  \source{petersen:page-0067,petersen:page-0068}
\end{proposition}
\begin{proof}
  \uses{prop:pet-ch2-hessian-euclidean-agreement}
  Write $\nabla f=\sum_i(\partial_if)\partial_i$. On the Cartesian coordinate frame,
  $L_{\partial_i}\partial_j=[\partial_i,\partial_j]=0$
  (\cref{prop:pet-ch2-lie-derivative-bracket}), so the only contribution to
  $L_{\nabla f}(dx^i)$ comes from the varying coefficient $\partial_if$ of
  $\nabla f=(\partial_if)\partial_i$, via \cref{prop:pet-ch2-lie-derivative-scaling}:
  \[
    L_{\nabla f}(\delta_{ij}dx^idx^j)
    = \sum_iL_{(\partial_if)\partial_i}(dx^i)dx^i
    = \sum_i\Big((L_{\partial_if}dx^i)dx^i+dx^i(L_{\partial_if}dx^i)\Big)
  \]
  \[
    = \sum_i\Big(\partial_if\,(L_{\partial_i}dx^i)dx^i+dx^i(\partial_if\,dx^i)\Big)
    = \sum_id(\partial_if)dx^i+\sum_id(\partial_if)dx^i
    = 2\sum_i\partial_{ij}f\,dx^idx^j,
  \]
  using that $L_{\partial_i}dx^i=0$ (dual to the vanishing bracket) so the
  $(\partial_if)$-derivative term dominates, that $d(\partial_if)=\partial_{ij}f\,dx^j$,
  and collecting the two symmetric contributions. Hence
  $\operatorname{Hess}f=\frac12L_{\nabla f}g=\sum_i\partial_{ij}f\,dx^idx^j$.
\end{proof}

\subsection{Lie Groups}

\begin{lemma}[Lem. 2.1.7 — differential of $\operatorname{Ad}$]
  \label{lem:pet-ch2-ad-differential}
  \lean{PetersenLib.ad_eq_differential_of_Ad}
  For a Lie group $G$ with inner automorphism $\operatorname{Ad}_h:x\mapsto hxh^{-1}$
  and its differential $\operatorname{Ad}_h:\mathfrak g\to\mathfrak g$ at $e$, the
  differential of $h\mapsto\operatorname{Ad}_h$ is $U\mapsto\operatorname{ad}_U(X)=[U,X]$.
  \source{petersen:page-0068}
\end{lemma}
\begin{proof}
  \uses{lem:pet-ch2-ad-differential}
  Write $\operatorname{Ad}_h(x)=R_{h^{-1}}L_h(x)$, so
  $\operatorname{Ad}_h=DR_{h^{-1}}DL_h$ at $e$. Let $F^t$ be the flow of the
  left-invariant field $U$; then $F^t(x)=xF^t(e)=L_x(F^t(e))$, since both sides
  agree at $t=0$ and have $U$ as tangent everywhere (left-invariance), and
  consequently $DF^t=DR_{F^t(e)}$. Hence, for $X\in\mathfrak g$,
  \[
    \operatorname{ad}_U(X)|_e
    = \frac{d}{dt}\Big(DR_{F^t(e)}DL_{F^t(e)}(X|_e)\Big)\Big|_{t=0}
    = \frac{d}{dt}\Big(DR_{F^t(e)}\big(X|_{F^t(e)}\big)\Big)\Big|_{t=0}
    = \frac{d}{dt}\Big(DF^{-t}\big(X|_{F^t(e)}\big)\Big)\Big|_{t=0}
  \]
  \[
    = L_UX = [U,X],
  \]
  using $\operatorname{Ad}_h(X)=DR_{h^{-1}}DL_h(X)$ and
  \cref{def:pet-ch2-lie-derivative-vector-field,prop:pet-ch2-lie-derivative-bracket}.
\end{proof}

\begin{lemma}[Lem. 2.1.8 — bracket on $\mathfrak{gl}(V)$]
  \label{lem:pet-ch2-gl-bracket-commutator}
  \lean{PetersenLib.gl_bracket_eq_commutator}
  \leanok
  \uses{lem:pet-ch2-ad-differential}
  Let $G=\mathrm{GL}(V)$. The Lie bracket on $\mathfrak{gl}(V)=T_I\mathrm{GL}(V)$
  of left-invariant vector fields is given by commutation of linear maps:
  $[X,Y]_{\mathfrak l}=XY-YX$.
  \source{petersen:page-0068,petersen:page-0069}
\end{lemma}
\begin{proof}
  \uses{lem:pet-ch2-gl-bracket-commutator}
  Since $x\mapsto hxh^{-1}$ is linear on $\mathrm{Hom}(V,V)$,
  $\operatorname{Ad}_h(X)=hXh^{-1}$. The flow of $U\in\mathfrak{gl}(V)$ is
  $F^t(g)=g(I+tU+o(t))$, so by \cref{lem:pet-ch2-ad-differential},
  \[
    [U,X] = \frac{d}{dt}\Big(F^t(I)\,X\,F^{-t}(I)\Big)\Big|_{t=0}
    = \frac{d}{dt}\Big((I+tU+o(t))X(I-tU+o(t))\Big)\Big|_{t=0}
    = \frac{d}{dt}\big(X+tUX-tXU+o(t)\big)\Big|_{t=0}
    = UX-XU.
  \]
\end{proof}

\section{Connections}

\subsection{Covariant Differentiation}

\begin{definition}[covariant derivative on $\mathbb{R}^n$]
  \label{def:pet-ch2-covariant-derivative-euclidean}
  \lean{PetersenLib.covariantDerivativeEuclidean}
  \leanok
  \uses{def:pet-ch2-directional-derivative}
  Writing $X=X^i\partial_i$ in Cartesian coordinates, the \emph{covariant
  derivative} of $X$ in the direction $Y$ is
  $\nabla_YX=(\nabla_YX^i)\partial_i=d(X^i)(Y)\partial_i$, so a vector field with
  constant coefficients has zero covariant derivative. This formula is not
  invariant under change of coordinates: e.g.\ the polar-coordinate frame
  fields $\partial_r=\frac1r(x\partial_x+y\partial_y)$, $\partial_\theta=-y\partial_x+x\partial_y$
  are not constant, yet ought to have a coordinate-independent covariant
  derivative.
  \source{petersen:page-0069}
\end{definition}

\begin{definition}[dual $1$-form $\theta_X$]
  \label{def:pet-ch2-dual-one-form}
  \lean{PetersenLib.dualOneForm}
  \leanok
  For a vector field $X$ on $(M,g)$, the dual $1$-form is
  $\theta_X:=i_Xg$, i.e.\ $\theta_X(Y)=g(X,Y)$. $X$ is locally a gradient field
  iff $d\theta_X=0$; in general $d\theta_X$ measures the extent to which $X$
  fails to be a gradient field.
  \source{petersen:page-0069}
\end{definition}

\begin{definition}[Killing field]
  \label{def:pet-ch2-killing-field}
  \lean{PetersenLib.IsKillingField}
  \leanok
  \uses{def:pet-ch2-lie-derivative-tensor}
  A vector field $X$ on $(M,g)$ is a \emph{Killing field} if $L_Xg=0$;
  equivalently, if $F^t$ is the local flow of $X$, the $F^t$ are local
  isometries.
  \source{petersen:page-0069}
\end{definition}

\begin{remark}[Killing fields on $\mathbb{R}^n$]
  \label{rem:pet-ch2-killing-field-euclidean}
  \lean{PetersenLib.killingField_euclidean_characterization}
  \uses{def:pet-ch2-killing-field}
  If $X=X^i\partial_i$ on $\mathbb{R}^n$, $X$ is Killing iff $\partial_iX^j+\partial_jX^i=0$ for
  all $i,j$. This forces $\partial_j\partial_iX^i=0$ for all $i,j$ (differentiating
  the skew-symmetry relation and permuting indices cyclically shows the
  third-order-derivative expression equals its own negative), so
  $X^i=\alpha^i_jx^j+\beta^i$ with $\alpha^i_j=\partial_jX^i=-\partial_iX^j=-\alpha^j_i$
  skew. In particular $\partial_\theta$ is a Killing field on $\mathbb{R}^2$, matching that
  its flow is rotation by an orthogonal matrix; more generally the flow of a
  Killing field on $\mathbb{R}^n$ is $F^t(x)=\exp(At)x+t\beta$ with $A=[\alpha^i_j]$
  skew-symmetric. A vector field on $\mathbb{R}^n$ is constant iff it is both Killing
  and a gradient field.
  \source{petersen:page-0070}
\end{remark}

\begin{proposition}[Prop. 2.2.1 — implicit formula on $\mathbb{R}^n$]
  \label{prop:pet-ch2-covariant-derivative-implicit-euclidean}
  \lean{PetersenLib.covariantDerivative_implicit_euclidean}
  \leanok
  \uses{def:pet-ch2-covariant-derivative-euclidean,def:pet-ch2-dual-one-form}
  The covariant derivative on $\mathbb{R}^n$ satisfies
  $2g(\nabla_YX,Z)=(L_Xg)(Y,Z)+(d\theta_X)(Y,Z)$.
  \source{petersen:page-0070,petersen:page-0071}
\end{proposition}
\begin{proof}
  \uses{prop:pet-ch2-covariant-derivative-implicit-euclidean}
  \leanok
  Both sides are tensorial in $Y,Z$, so it suffices to check on Cartesian
  coordinate fields $\partial_k,\partial_\ell$. Write $X=a^i\partial_i$. Since
  $L_{\partial_i}\partial_k=[\partial_i,\partial_k]=0$
  (\cref{prop:pet-ch2-lie-derivative-bracket}) and, by
  \cref{prop:pet-ch2-lie-derivative-scaling}, $L_{a^i\partial_i}\partial_k=-(\partial_ka^i)\partial_i$,
  \[
    (L_Xg)(\partial_k,\partial_\ell)+(d\theta_X)(\partial_k,\partial_\ell)
    = D_X\delta_{k\ell}-g(L_X\partial_k,\partial_\ell)-g(\partial_k,L_X\partial_\ell)
    +\partial_kg(X,\partial_\ell)-\partial_\ell g(X,\partial_k)-g(X,[\partial_k,\partial_\ell])
  \]
  \[
    = -g\big(-(\partial_ka^i)\partial_i,\partial_\ell\big)-g\big(\partial_k,-(\partial_\ell a^i)\partial_i\big)
    +\partial_ka^\ell-\partial_\ell a^k
    = \partial_ka^\ell+\partial_\ell a^k+\partial_ka^\ell-\partial_\ell a^k
    = 2\partial_ka^\ell
  \]
  \[
    = 2g\big((\partial_ka^\ell)\partial_\ell,\partial_\ell\big)
    = 2g(\nabla_{\partial_k}X,\partial_\ell),
  \]
  using $D_X\delta_{k\ell}=0$, $g(X,\partial_\ell)=a^\ell$, and
  \cref{def:pet-ch2-covariant-derivative-euclidean}.
\end{proof}

\begin{definition}[affine connection]
  \label{def:pet-ch2-affine-connection}
  \lean{PetersenLib.AffineConnection}
  \leanok
  On a manifold $M$, an \emph{affine connection} is an assignment
  $X\mapsto\nabla X$, $\nabla:TM\times\mathfrak X(M)\to TM$, such that: (1)
  $v\mapsto\nabla_vX$ is a $(1,1)$-tensor, i.e.\ linear,
  $\nabla_{\alpha v+\beta w}X=\alpha\nabla_vX+\beta\nabla_wX$; and (2)
  $X\mapsto\nabla_YX$ is a derivation, $\nabla_Y(X_1+X_2)=\nabla_YX_1+\nabla_YX_2$
  and $\nabla_Y(fX)=(D_Yf)X+f\nabla_YX$ for $f:M\to\mathbb{R}$.
  \source{petersen:page-0071,petersen:page-0073}
\end{definition}

\begin{definition}[Riemannian connection]
  \label{def:pet-ch2-riemannian-connection}
  \lean{PetersenLib.RiemannianConnection}
  \leanok
  \uses{def:pet-ch2-affine-connection}
  If $(M,g)$ is Riemannian, an affine connection $\nabla$
  (\cref{def:pet-ch2-affine-connection}) is a \emph{Riemannian connection} if
  in addition: (3) it is torsion free, $\nabla_XY-\nabla_YX=[X,Y]$; and (4) it
  is metric, $D_Zg(X,Y)=g(\nabla_ZX,Y)+g(X,\nabla_ZY)$.
  \source{petersen:page-0071,petersen:page-0073}
\end{definition}

\begin{remark}[Koszul's formula]
  \label{rem:pet-ch2-koszul-formula}
  \lean{PetersenLib.koszulFormula}
  \leanok
  \uses{prop:pet-ch2-covariant-derivative-implicit-euclidean,def:pet-ch2-dual-one-form,def:pet-ch2-lie-derivative-tensor}
  On a general Riemannian manifold, the expression
  $(L_Xg)(Y,Z)+(d\theta_X)(Y,Z)$ generalizing
  \cref{prop:pet-ch2-covariant-derivative-implicit-euclidean} expands, by the
  formulas of \cref{prop:pet-ch2-lie-derivative-tensor-formula} for $L_Xg$ and
  \cref{def:pet-ch2-exterior-derivative-lie-formula} for $d\theta_X$, to
  \[
    (L_Xg)(Y,Z)+(d\theta_X)(Y,Z)
    = D_Xg(Y,Z)-g([X,Y],Z)-g(Y,[X,Z])
  \]
  \[
    +D_Y\theta_X(Z)-D_Z\theta_X(Y)-\theta_X([Y,Z])
    = D_Xg(Y,Z)+D_Yg(X,Z)-D_Zg(X,Y)
    -g([X,Y],Z)-g([Y,Z],X)+g([Z,X],Y).
  \]
  \source{petersen:page-0072}
\end{remark}


\begin{lemma}[extension of a tangent vector to a vector field]
  \label{lem:pet-ch2-extend-tangent-vector}
  \lean{PetersenLib.exists_smoothVectorField_eq}
  \leanok
  On a smooth ($\sigma$-compact, Hausdorff) manifold, every tangent vector
  $v\in T_pM$ is the value at $p$ of a globally defined smooth vector field.
\end{lemma}
\begin{proof}
  \leanok
  Extend the constant coefficient vector field through a chart at $p$ and
  multiply by a smooth bump function supported in the chart and equal to $1$
  near $p$; the resulting field is globally smooth and takes the value $v$
  at $p$.
\end{proof}

\begin{lemma}[Riemannian connections satisfy Koszul's formula]
  \label{lem:pet-ch2-riemannian-connection-koszul}
  \lean{PetersenLib.RiemannianConnection.koszul}
  \leanok
  \uses{def:pet-ch2-riemannian-connection,rem:pet-ch2-koszul-formula}
  Every Riemannian connection $\nabla$ on $(M,g)$ satisfies
  $2g(\nabla_YX,Z)=D_Xg(Y,Z)+D_Yg(X,Z)-D_Zg(X,Y)
  -g([X,Y],Z)-g([Y,Z],X)+g([Z,X],Y)$.
\end{lemma}
\begin{proof}
  \leanok
  Expand the three derivative terms by the metric property and the three
  bracket terms by torsion-freeness; all terms cancel except
  $2g(\nabla_YX,Z)$, using the symmetry of $g$.
\end{proof}

\begin{theorem}[Thm. 2.2.2 — the Fundamental Theorem of Riemannian Geometry]
  \label{thm:pet-ch2-fundamental-theorem}
  \lean{PetersenLib.fundamentalTheoremRiemannianGeometry}
  \leanok
  \uses{def:pet-ch2-riemannian-connection,rem:pet-ch2-koszul-formula,lem:pet-ch2-extend-tangent-vector,lem:pet-ch2-riemannian-connection-koszul}
  On $(M,g)$ there is exactly one Riemannian connection $\nabla$
  (\cref{def:pet-ch2-riemannian-connection}), given implicitly by Koszul's
  formula (\cref{rem:pet-ch2-koszul-formula}):
  \[
    2g(\nabla_YX,Z)=D_Xg(Y,Z)+D_Yg(X,Z)-D_Zg(X,Y)
    -g([X,Y],Z)-g([Y,Z],X)+g([Z,X],Y).
  \]
  \source{petersen:page-0071,petersen:page-0072,petersen:page-0073}
\end{theorem}
\begin{proof}
  \uses{thm:pet-ch2-fundamental-theorem}
  \leanok
  \emph{Existence, properties (1)–(2).} Define $\nabla_YX$ by
  $2g(\nabla_YX,Z):=(L_Xg)(Y,Z)+(d\theta_X)(Y,Z)$, tensorial and linear in
  $Y,Z$ by construction (\cref{rem:pet-ch2-koszul-formula}), hence linear in
  $Y$ (property (1)). For the derivation property in $X$, since
  $L_{fX}g+d\theta_{fX}=fL_Xg+df\wedge\theta_X+\theta_X\cdot df+d(f\theta_X)
  =fL_Xg+df\cdot\theta_X+\theta_X\cdot df+df\wedge\theta_X+fd\theta_X
  =f(L_Xg+d\theta_X)+2df\cdot\theta_X$, evaluating on $(Y,Z)$ gives
  $2g(\nabla_Y(fX),Z)=f\cdot2g(\nabla_YX,Z)+2df(Y)g(X,Z)
  =2g\big(f\nabla_YX+(D_Yf)X,Z\big)$, so
  $\nabla_Y(fX)=f\nabla_YX+(D_Yf)X$ (property (2)).

  \emph{Torsion-free (3).} By \cref{rem:pet-ch2-koszul-formula} applied twice
  and antisymmetrizing in $X,Y$, all the symmetric $D$-terms cancel and the
  bracket terms combine as
  \[
    2g(\nabla_XY-\nabla_YX,Z)
    = D_Yg(X,Z)+D_Yg(Z,X)-D_Zg(Y,X)-g([Y,X],Z)-g([X,Z],Y)+g([Z,X],Y)
  \]
  \[
    -D_Xg(Y,Z)-D_Yg(Z,X)+D_Zg(X,Y)+g([X,Y],Z)+g([Y,Z],X)-g([Z,Y],X)
    = 2g([X,Y],Z),
  \]
  so $\nabla_XY-\nabla_YX=[X,Y]$.

  \emph{Metric (4).} Symmetrizing \cref{rem:pet-ch2-koszul-formula} in $X,Y$
  instead,
  \[
    2g(\nabla_ZX,Y)+2g(X,\nabla_ZY)
    = D_Xg(Z,Y)+D_Zg(Y,X)-D_Yg(X,Z)-g([X,Z],Y)-g([Z,Y],X)+g([Y,Z],X)
  \]
  \[
    +D_Yg(Z,X)+D_Zg(X,Y)-D_Yg(X,Z)-g([Y,Z],X)-g([Z,X],Y)+g([X,Y],Z)
    = 2D_Zg(X,Y),
  \]
  so $D_Zg(X,Y)=g(\nabla_ZX,Y)+g(X,\nabla_ZY)$.

  \emph{Uniqueness.} If $\tilde\nabla$ also satisfies (1)–(4), then applying
  metricity and torsion-freeness of $\tilde\nabla$ to expand
  $D_Xg(Y,Z)+D_Yg(Z,X)-D_Zg(X,Y)$ term by term and substituting brackets
  $[X,Y]=\tilde\nabla_XY-\tilde\nabla_YX$ etc.\ into
  \cref{rem:pet-ch2-koszul-formula} produces, after cancellation of all
  cross-terms, $2g(\nabla_YX,Z)=2g(\tilde\nabla_YX,Z)$ for all $Z$, i.e.\
  $\nabla_YX=\tilde\nabla_YX$; explicitly, substituting the metricity identity
  four times and the torsion identity twice into the right-hand side of
  Koszul's formula collapses it to $2g(\tilde\nabla_YX,Z)$, matching the
  left-hand side that defines $\nabla$.
\end{proof}

\begin{lemma}[Lem. 2.2.3 — locality on open sets]
  \label{lem:pet-ch2-locality-open-set}
  \lean{PetersenLib.connection_local_openSet}
  \leanok
  \uses{def:pet-ch2-affine-connection}
  Let $\nabla$ be an affine connection on $M$, $p\in M$, $v\in T_pM$, and
  $X,Y$ vector fields with $X=Y$ on a neighborhood $U\ni p$. Then
  $\nabla_vX=\nabla_vY$.
  \source{petersen:page-0073,petersen:page-0074}
\end{lemma}
\begin{proof}
  \uses{lem:pet-ch2-locality-open-set}
  \leanok
  Choose $\lambda:M\to\mathbb{R}$ with $\lambda\equiv0$ off $U$ and $\lambda\equiv1$
  near $p$, so $\lambda X=\lambda Y$ on $M$. By the derivation property,
  $\nabla_v(\lambda X)=\lambda(p)\nabla_vX+d\lambda(v)X(p)=\nabla_vX$ since
  $d\lambda|_p=0$ and $\lambda(p)=1$; likewise $\nabla_v(\lambda Y)=\nabla_vY$.
  Since $\lambda X=\lambda Y$ as global fields,
  $\nabla_vX=\nabla_v(\lambda X)=\nabla_v(\lambda Y)=\nabla_vY$.
\end{proof}

\begin{lemma}[Lem. 2.2.4 — locality along a curve]
  \label{lem:pet-ch2-locality-along-curve}
  \lean{PetersenLib.connection_local_alongCurve}
  \uses{def:pet-ch2-affine-connection}
  Let $\nabla$ be an affine connection, $X$ a vector field, and $c:I\to M$
  smooth with $\dot c(0)=v\in T_pM$. If $X\circ c=Y\circ c$ for another vector
  field $Y$, then $\nabla_vX=\nabla_vY$.
  \source{petersen:page-0074}
\end{lemma}
\begin{proof}
  \uses{lem:pet-ch2-locality-along-curve}
  Choose a frame $E_1,\dots,E_n$ near $p$ and write $X=X^iE_i$, $Y=Y^iE_i$.
  Since $X\circ c=Y\circ c$, the coefficient functions satisfy
  $X^i\circ c=Y^i\circ c$. Then, by the derivation and linearity properties
  of $\nabla$,
  \[
    \nabla_vY = \nabla_v(Y^iE_i) = Y^i(p)\nabla_vE_i+E_i(p)\,dY^i(v)
    = X^i(p)\nabla_vE_i+E_i(p)\,dX^i(v) = \nabla_vX,
  \]
  using $Y^i(p)=X^i(p)$ (from $c(0)=p$) and $dY^i(v)=dX^i(v)$ (since
  $Y^i\circ c=X^i\circ c$ as functions of the curve parameter, so their
  derivatives at $0$ along $\dot c(0)=v$ agree).
\end{proof}

\begin{definition}[normal coordinates and normal frame at $p$]
  \label{def:pet-ch2-normal-coordinates-frame}
  \lean{PetersenLib.isNormalCoordinatesAt,PetersenLib.isNormalFrameAt}
  \uses{def:pet-ch2-riemannian-connection}
  A coordinate system is \emph{normal} at $p$ if $g_{ij}|_p=\delta_{ij}$ and
  $\partial_kg_{ij}|_p=0$. An orthonormal frame $E_1,\dots,E_n$ is \emph{normal}
  at $p$ if $(\nabla_vE_i)|_p=0$ for all $i$ and $v\in T_pM$. Such coordinates
  and frames always exist at every point.
  \source{petersen:page-0074}
\end{definition}

\subsection{Covariant Derivatives of Tensors}

\begin{definition}[covariant derivative of an $(s,t)$-tensor field]
  \label{def:pet-ch2-covariant-derivative-tensor}
  \lean{PetersenLib.covariantDerivativeTensor}
  \leanok
  \uses{def:pet-ch2-riemannian-connection}
  For an $(s,t)$-tensor field $S$, the \emph{covariant derivative} $\nabla S$
  is the $(s,t+1)$-tensor field determined by requiring the Leibniz rule
  $\nabla_X(S(Y))=(\nabla_XS)(Y)+S(\nabla_XY)$: for $s=0,1$,
  \[
    \nabla S(X,Y_1,\dots,Y_r)=(\nabla_XS)(Y_1,\dots,Y_r)
    =\nabla_X\big(S(Y_1,\dots,Y_r)\big)-\sum_{i=1}^rS(Y_1,\dots,\nabla_XY_i,\dots,Y_r),
  \]
  where $\nabla_X$ acts as $D_X$ on functions and as covariant differentiation
  on vector fields; equivalently $\nabla_XS=[\nabla_X,S]$ when $S$ is
  $(1,1)$. For general $s$, extend by the product rule
  $\nabla_X(X_1\otimes X_2)=(\nabla_XX_1)\otimes X_2+X_1\otimes(\nabla_XX_2)$.
  \source{petersen:page-0075}
\end{definition}

\begin{definition}[parallel tensor]
  \label{def:pet-ch2-parallel-tensor}
  \lean{PetersenLib.IsParallel}
  \leanok
  \uses{def:pet-ch2-covariant-derivative-tensor}
  A tensor $S$ is \emph{parallel} if $\nabla S=0$. On Euclidean space, a
  tensor written in Cartesian coordinates is parallel iff it has constant
  coefficients.
  \source{petersen:page-0075}
\end{definition}

\begin{proposition}[Prop. 2.2.5 — the metric and volume form are parallel]
  \label{prop:pet-ch2-metric-volume-parallel}
  \lean{PetersenLib.metric_and_volume_parallel}
  \leanok
  \uses{def:pet-ch2-parallel-tensor}
  On a Riemannian $n$-manifold, $\nabla g=0$ and $\nabla\mathrm{vol}=0$.
  \source{petersen:page-0076}
\end{proposition}
\begin{proof}
  \uses{prop:pet-ch2-metric-volume-parallel}
  \leanok
  By the metric property (4) of \cref{def:pet-ch2-riemannian-connection},
  $(\nabla g)(X,Y_1,Y_2)=\nabla_X(g(Y_1,Y_2))-g(\nabla_XY_1,Y_2)-g(Y_1,\nabla_XY_2)=0$.
  For the volume form, evaluate on an orthonormal frame $E_1,\dots,E_n$:
  \[
    (\nabla_X\mathrm{vol})(E_1,\dots,E_n)
    = \nabla_X\big(\mathrm{vol}(E_1,\dots,E_n)\big)-\sum_i\mathrm{vol}(E_1,\dots,\nabla_XE_i,\dots,E_n)
    = -\sum_ig(E_i,\nabla_XE_i)
  \]
  \[
    = -\frac12\sum_iD_X\big(g(E_i,E_i)\big) = 0,
  \]
  using that $\mathrm{vol}(E_1,\dots,E_n)\equiv1$ and $g(E_i,E_i)\equiv1$ are
  constant, and again that replacing one frame slot projects onto that
  direction.
\end{proof}

\begin{proposition}[Prop. 2.2.6 — Hessian via the covariant derivative]
  \label{prop:pet-ch2-hessian-covariant-derivative}
  \lean{PetersenLib.hessian_via_covariantDerivative}
  \leanok
  \uses{def:pet-ch2-covariant-derivative-tensor,def:pet-ch2-hessian-lie-derivative}
  For $f:(M,g)\to\mathbb{R}$, $(\nabla df)(X,Y)=g(\nabla_X\nabla f,Y)=\operatorname{Hess}f(X,Y)$.
  \source{petersen:page-0076,petersen:page-0077}
\end{proposition}
\begin{proof}
  \uses{prop:pet-ch2-hessian-covariant-derivative}
  \leanok
  By \cref{def:pet-ch2-covariant-derivative-tensor},
  $(\nabla df)(X,Y)=(\nabla_Xdf)(Y)=D_XD_Yf-df(\nabla_XY)=D_XD_Yf-D_{\nabla_XY}f$,
  so
  \[
    (\nabla_Xdf)(Y)-(\nabla_Ydf)(X)=[D_X,D_Y]f-D_{[X,Y]}f=0
  \]
  since $\nabla$ is torsion free (\cref{def:pet-ch2-riemannian-connection}(3)),
  showing $\nabla df$ is symmetric. Next, using the metric property,
  \[
    (\nabla_Xdf)(Y)=D_Xg(\nabla f,Y)-g(\nabla f,\nabla_XY)=g(\nabla_X\nabla f,Y).
  \]
  On the other hand, expanding via the metric property and the torsion-free
  bracket identity as in \cref{prop:pet-ch2-covariant-derivative-implicit-euclidean},
  \[
    g(\nabla_X\nabla f,Y)
    = \tfrac12g(\nabla_X\nabla f,Y)+\tfrac12g(X,\nabla_Y\nabla f)
    + \tfrac12D_{\nabla f}g(X,Y)-\tfrac12g([\nabla f,X],Y)-\tfrac12g(X,[\nabla f,Y])
  \]
  \[
    = \tfrac12(L_{\nabla f}g)(X,Y),
  \]
  matching $\operatorname{Hess}f(X,Y)=\frac12(L_{\nabla f}g)(X,Y)$ from
  \cref{def:pet-ch2-hessian-lie-derivative}.
\end{proof}

\subsection{The Adjoint of the Covariant Derivative}

\begin{definition}[adjoint of the covariant derivative]
  \label{def:pet-ch2-covariant-derivative-adjoint}
  \lean{PetersenLib.covariantDerivativeAdjoint}
  \uses{def:pet-ch2-covariant-derivative-tensor}
  For an $(s,t)$-tensor $S$ with $t>0$, and an orthonormal frame
  $E_1,\dots,E_n$, the \emph{adjoint} $\nabla^*S$ is the $(s,t-1)$-tensor
  \[
    (\nabla^*S)(X_2,\dots,X_t)=-\sum_i(\nabla_{E_i}S)(E_i,X_2,\dots,X_t).
  \]
  \source{petersen:page-0077}
\end{definition}

\begin{proposition}[Prop. 2.2.7 — divergence as adjoint of $\theta_X$]
  \label{prop:pet-ch2-divergence-adjoint-formula}
  \lean{PetersenLib.divergence_eq_neg_adjoint_dualForm}
  \uses{def:pet-ch2-covariant-derivative-adjoint,def:pet-ch2-divergence-lie-derivative,def:pet-ch2-dual-one-form}
  If $X$ is a vector field with dual $1$-form $\theta_X$, then
  $\operatorname{div}X=-\nabla^*\theta_X$.
  \source{petersen:page-0077}
\end{proposition}
\begin{proof}
  \uses{prop:pet-ch2-divergence-adjoint-formula}
  Select an orthonormal frame $E_i$. By
  \cref{def:pet-ch2-covariant-derivative-adjoint,def:pet-ch2-covariant-derivative-tensor},
  \[
    -\nabla^*\theta_X = \sum_i(\nabla_{E_i}\theta_X)(E_i)
    = \sum_iD_{E_i}g(X,E_i)-\sum_ig(X,\nabla_{E_i}E_i)
    = \sum_ig(\nabla_{E_i}X,E_i)
  \]
  \[
    = \sum_i\tfrac12(L_Xg)(E_i,E_i) = \operatorname{div}X,
  \]
  where the third equality follows exactly as in
  \cref{prop:pet-ch2-hessian-covariant-derivative}, pairing the connection
  with the metric, and the last equality is
  \cref{prop:pet-ch2-divergence-trace-formula}.
\end{proof}

\begin{proposition}[Prop. 2.2.8 — $\nabla^*$ is the $L^2$-adjoint of $\nabla$]
  \label{prop:pet-ch2-adjoint-integral-identity}
  \lean{PetersenLib.adjoint_L2_identity}
  \uses{def:pet-ch2-covariant-derivative-adjoint}
  If $S$ is a compactly supported $(s,t)$-tensor and $T$ a compactly supported
  $(s,t+1)$-tensor, then $\int g(\nabla S,T)\,\mathrm{vol}=\int g(S,\nabla^*T)\,\mathrm{vol}$.
  \source{petersen:page-0078}
\end{proposition}
\begin{proof}
  \uses{prop:pet-ch2-adjoint-integral-identity}
  For notational simplicity take $s=t=1$; the general case is identical with
  more indices. Define the $1$-form $\omega(X)=g(i_XT,S)$. At $p$, choose an
  orthonormal frame $E_i$ normal at $p$ (\cref{def:pet-ch2-normal-coordinates-frame}),
  so $\nabla_vE_i|_p=0$ for all $v$. Then at $p$,
  \[
    -\nabla^*\omega = \sum_i(\nabla_{E_i}\omega)(E_i)
    = \sum_i\nabla_{E_i}g\big(T(E_i,E_j),S(E_j)\big)
    = \sum_ig\big(\nabla_{E_i}T(E_i,E_j),S(E_j)\big)+g\big(T(E_i,E_j),\nabla_{E_i}S(E_j)\big)
  \]
  \[
    = -g(\nabla^*T,S)+g(T,\nabla S).
  \]
  Since this pointwise identity holds at every $p$ (using a frame normal at
  that point), integrating and applying the divergence theorem
  $\int\operatorname{div}(\theta^\sharp)\,\mathrm{vol}=\int d\,i_{\theta^\sharp}\mathrm{vol}=0$
  (\cref{rem:pet-ch2-lie-derivative-volume-exact}, Stokes' theorem, valid since
  $\omega$ has compact support as $S,T$ do) to $-\nabla^*\omega=\operatorname{div}(\omega^\sharp)$
  gives $\int(-g(\nabla^*T,S)+g(T,\nabla S))\,\mathrm{vol}=0$, i.e.\
  $\int g(\nabla S,T)\,\mathrm{vol}=\int g(S,\nabla^*T)\,\mathrm{vol}$.
\end{proof}

\subsection{Exterior Derivatives}

\begin{definition}[exterior derivative via the covariant derivative]
  \label{def:pet-ch2-exterior-derivative-covariant}
  \lean{PetersenLib.exteriorDerivative_covariantFormula}
  \uses{def:pet-ch2-covariant-derivative-tensor}
  For a $k$-form $\omega$, $(d\omega)(X_0,\dots,X_k)=\sum_i(-1)^i(\nabla_{X_i}\omega)(X_0,\dots,\hat X_i,\dots,X_k)$.
  More generally, for a $(1,k)$-tensor $T$ skew-symmetric in its $k$
  covariant slots, define the $(1,k+1)$-tensor
  $(d^\nabla T)(X_0,\dots,X_k)=\sum_i(-1)^i(\nabla_{X_i}T)(X_0,\dots,\hat X_i,\dots,X_k)$;
  for $k=0$, $T=Y$ a vector field, $(d^\nabla Y)(X)=\nabla_XY$; for $k=1$,
  \[
    (d^\nabla T)(X,Y)=(\nabla_XT)(Y)-(\nabla_YT)(X)
    =\nabla_X(T(Y))-\nabla_Y(T(X))-T([X,Y]).
  \]
  \source{petersen:page-0078,petersen:page-0079}
\end{definition}

\subsection{The Second Covariant Derivative}

\begin{definition}[second covariant derivative]
  \label{def:pet-ch2-second-covariant-derivative}
  \lean{PetersenLib.secondCovariantDerivative}
  \leanok
  \uses{def:pet-ch2-covariant-derivative-tensor}
  For an $(s,t)$-tensor field $S$, the \emph{second covariant derivative}
  $\nabla^2S$ is the $(s,t+2)$-tensor field
  \[
    (\nabla^2_{X_1,X_2}S)(Y_1,\dots,Y_t)
    = (\nabla_{X_1}(\nabla S))(X_2,Y_1,\dots,Y_t)
    = (\nabla_{X_1}(\nabla_{X_2}S))(Y_1,\dots,Y_t)-(\nabla_{\nabla_{X_1}X_2}S)(Y_1,\dots,Y_t).
  \]
  \source{petersen:page-0079}
\end{definition}

\begin{remark}[Hessian as second covariant derivative; Laplacian as trace]
  \label{rem:pet-ch2-hessian-second-covariant-derivative}
  \lean{PetersenLib.hessian_eq_secondCovariantDerivative}
  \uses{def:pet-ch2-second-covariant-derivative,def:pet-ch2-laplacian,def:pet-ch2-covariant-derivative-adjoint}
  For $f:M\to\mathbb{R}$, $\nabla^2_{X,Y}f=\nabla_X\nabla_Yf-\nabla_{\nabla_XY}f
  =\nabla_Xdf(Y)-df(\nabla_XY)=(\nabla_Xdf)(Y)=\operatorname{Hess}f(X,Y)$, and
  $\nabla^2f$ is symmetric in $X,Y$ (unlike the second covariant derivative of
  a general tensor, whose failure of symmetry is central to curvature).
  Consequently $\Delta f=-\nabla^*\nabla f=\sum_i\nabla^2_{E_i,E_i}f$.
  \source{petersen:page-0079}
\end{remark}

\subsection{The Lie Derivative of the Covariant Derivative}

\begin{definition}[Lie derivative of the connection]
  \label{def:pet-ch2-lie-derivative-connection}
  \lean{PetersenLib.lieDerivativeConnection}
  \leanok
  \uses{def:pet-ch2-lie-derivative-vector-field,def:pet-ch2-riemannian-connection}
  \[
    (L_X\nabla)_UV := (L_X\nabla)(U,V)
    = L_X(\nabla_UV)-\nabla_{L_XU}V-\nabla_UL_XV
    = [X,\nabla_UV]-\nabla_{[X,U]}V-\nabla_U[X,V].
  \]
  \source{petersen:page-0079,petersen:page-0080}
\end{definition}

\begin{remark}[symmetry and tensoriality of $L_X\nabla$]
  \label{rem:pet-ch2-lie-derivative-connection-symmetric}
  \lean{PetersenLib.lieDerivativeConnection_symmetric}
  \leanok
  \uses{def:pet-ch2-lie-derivative-connection}
  $(L_X\nabla)(U,V)$ is tensorial in $U$ (since $\nabla_UV$ is tensorial in
  $U$), and it is symmetric in $U,V$; hence it is tensorial in both
  variables.
  \source{petersen:page-0080}
\end{remark}
\begin{proof}
  \uses{rem:pet-ch2-lie-derivative-connection-symmetric}
  \leanok
  Since $[U,V]=\nabla_UV-\nabla_VU$ (torsion-freeness,
  \cref{def:pet-ch2-riemannian-connection}(3)),
  \[
    (L_X\nabla)(U,V)-(L_X\nabla)(V,U)
    = L_X\big(\nabla_UV-\nabla_VU\big)-\nabla_{[X,U]}V+\nabla_{[X,V]}U-\nabla_U[X,V]+\nabla_V[X,U]
  \]
  \[
    = L_X([U,V])-\nabla_{[X,U]}V+\nabla_{[X,V]}U-\nabla_U[X,V]+\nabla_V[X,U]
    = L_X(L_UV) = 0,
  \]
  by the definition of $L_X\nabla$ applied to both orderings, the torsion
  identity, and the Generalized Jacobi Identity
  (\cref{prop:pet-ch2-generalized-jacobi-identity}) applied to $T=V$: expanding
  $(L_XL)_UV=L_X(L_UV)-L_{L_XU}V-L_U(L_XV)=0$ and comparing with the
  difference above (using $L_XU=[X,U]$, $L_XV=[X,V]$ by
  \cref{prop:pet-ch2-lie-derivative-bracket}) shows the difference equals
  $(L_XL)_UV=0$.
\end{proof}

\subsection{The Covariant Derivative of the Covariant Derivative}

\begin{definition}[covariant derivative of the connection]
  \label{def:pet-ch2-covariant-derivative-connection}
  \lean{PetersenLib.covariantDerivativeConnection}
  \leanok
  \uses{def:pet-ch2-covariant-derivative-tensor,def:pet-ch2-second-covariant-derivative}
  \[
    (\nabla_X\nabla)_YT := \nabla_X(\nabla_YT)-\nabla_{\nabla_XY}T-\nabla_Y(\nabla_XT).
  \]
  This is not tensorial in $X$. It relates to the second covariant derivative
  of $T$ (\cref{def:pet-ch2-second-covariant-derivative}) by
  $\nabla^2_{X,Y}T=(\nabla_X\nabla)_YT+\nabla_Y(\nabla_XT)$.
  \source{petersen:page-0080}
\end{definition}

\section{Natural Derivations}

\subsection{Endomorphisms as Derivations}

\begin{definition}[derivation on tensors]
  \label{def:pet-ch2-tensor-derivation}
  \lean{PetersenLib.IsTensorDerivation}
  \leanok
  A \emph{derivation} on tensors is a type-preserving linear map
  $T\mapsto DT$ that commutes with contractions and satisfies the product rule
  $D(T_1\otimes T_2)=(DT_1)\otimes T_2+T_1\otimes(DT_2)$.
  \source{petersen:page-0080}
\end{definition}

\begin{definition}[the natural $\mathrm{GL}(V)$-action on $T(V)$]
  \label{def:pet-ch2-gl-action-tensor-algebra}
  \lean{PetersenLib.glActionOnTensorAlgebra}
  \leanok
  For a finite-dimensional vector space $V$, let $T(V)$ be its tensor algebra,
  graded by $(s,t)$-tensors spanned by
  $v_1\otimes\cdots\otimes v_s\otimes\phi_1\otimes\cdots\otimes\phi_t$
  ($v_i\in V$, $\phi_j\in V^*$). The natural homomorphism
  $\mathrm{GL}(V)\to\mathrm{GL}(T(V))$ acts by $g\cdot\alpha=\alpha$ for
  $\alpha\in\mathbb{R}$ (a $(0,0)$-tensor); $g\cdot v=g(v)$ for $v\in V$;
  $g\cdot\phi=\phi\circ g^{-1}$ for $\phi\in V^*$; and, on general tensors,
  \[
    g\cdot(v_1\otimes\cdots\otimes v_s\otimes\phi_1\otimes\cdots\otimes\phi_t)
    = g(v_1)\otimes\cdots\otimes g(v_s)\otimes(\phi_1\circ g^{-1})\otimes\cdots\otimes(\phi_t\circ g^{-1}).
  \]
  \source{petersen:page-0080,petersen:page-0081}
\end{definition}

\begin{definition}[induced derivation of $L\in\mathrm{End}(V)$]
  \label{def:pet-ch2-endomorphism-derivation}
  \lean{PetersenLib.endomorphismDerivation}
  \leanok
  \uses{def:pet-ch2-gl-action-tensor-algebra,def:pet-ch2-tensor-derivation}
  Differentiating the $\mathrm{GL}(V)$-action of
  \cref{def:pet-ch2-gl-action-tensor-algebra} gives a linear map
  $\mathrm{End}(V)\to\mathrm{End}(T(V))$, $L\mapsto LT$: on vectors $Lv=L(v)$;
  on $1$-forms $L\phi=-\phi\circ L$; and on general tensors,
  \[
    L(v_1\otimes\cdots\otimes v_s\otimes\phi_1\otimes\cdots\otimes\phi_t)
    = \sum_{i=1}^sv_1\otimes\cdots\otimes L(v_i)\otimes\cdots\otimes v_s\otimes\phi_1\otimes\cdots\otimes\phi_t
  \]
  \[
    -\sum_{j=1}^tv_1\otimes\cdots\otimes v_s\otimes\phi_1\otimes\cdots\otimes(\phi_j\circ L)\otimes\cdots\otimes\phi_t.
  \]
  \source{petersen:page-0081}
\end{definition}

\begin{proposition}[Prop. 2.3.1 — Lie algebra homomorphism]
  \label{prop:pet-ch2-endomorphism-lie-algebra-hom}
  \lean{PetersenLib.endomorphismDerivation_lieAlgebraHom}
  \leanok
  \uses{def:pet-ch2-endomorphism-derivation}
  The linear map $\mathrm{End}(V)\to\mathrm{End}(T(V))$, $L\mapsto LT$, is a
  Lie algebra homomorphism that preserves symmetries of tensors.
  \source{petersen:page-0081}
\end{proposition}

\begin{proposition}[Prop. 2.3.2 — $(1,1)$-tensors act as derivations]
  \label{prop:pet-ch2-endomorphism-is-derivation}
  \lean{PetersenLib.endomorphismAction_isDerivation}
  \leanok
  \uses{def:pet-ch2-endomorphism-derivation,def:pet-ch2-tensor-derivation}
  Any $(1,1)$-tensor $L$ defines a derivation on tensors
  (\cref{def:pet-ch2-tensor-derivation}).
  \source{petersen:page-0081,petersen:page-0082}
\end{proposition}
\begin{proof}
  \uses{prop:pet-ch2-endomorphism-is-derivation}
  \leanok
  Linearity and the product rule are immediate from
  \cref{def:pet-ch2-endomorphism-derivation}. It remains to check that $L$
  commutes with contractions. In a local frame $X_i$ with coframe $\sigma^i$,
  write a $(1,1)$-tensor $T=T^i_jX_i\otimes\sigma^j$; its contraction is
  $\operatorname{tr}T=T^i_i$, a scalar, so $L(\operatorname{tr}T)=0$ (scalars
  are annihilated by $L$). On the other hand, writing $L(X_i)=L^k_iX_k$,
  \[
    L(T) = T^i_jL(X_i)\otimes\sigma^j-T^i_jX_i\otimes(\sigma^j\circ L)
    = T^i_jL^k_iX_k\otimes\sigma^j-T^i_jL^j_kX_i\otimes\sigma^k
  \]
  \[
    = T^i_jL^k_iX_k\otimes\sigma^j-T^k_jL^j_iX_i\otimes\sigma^j
    = \big(T^i_jL^k_i-T^k_jL^j_i\big)X_k\otimes\sigma^j,
  \]
  where the second sum was reindexed $(i,k)\mapsto(k,i)$. Contracting
  ($k=j$),
  \[
    \operatorname{tr}(LT) = T^i_kL^k_i-T^k_kL^k_i\big|_{j=k}
    = T^i_kL^k_i-T^i_kL^k_i = 0,
  \]
  matching $L(\operatorname{tr}T)=0$. The same computation, tracing over one
  fixed upper and one fixed lower index, applies to a general tensor
  $T^{j\cdots k}_{h\cdots i}$.
\end{proof}

\begin{definition}[inner product on the tensor algebra]
  \label{def:pet-ch2-tensor-algebra-inner-product}
  \lean{PetersenLib.tensorAlgebraInnerProduct}
  \leanok
  If $V$ has an inner product with orthonormal basis $e_1,\dots,e_n$ and dual
  basis $e^1,\dots,e^n$, the inner product on $T(V)$ declares
  $e_{i_1}\otimes\cdots\otimes e_{i_p}\otimes e^{j_1}\otimes\cdots\otimes e^{j_q}$
  an orthonormal basis of the $(p,q)$-tensors.
  \source{petersen:page-0082}
\end{definition}

\begin{proposition}[Prop. 2.3.3 — adjoints and skew-adjoint endomorphisms]
  \label{prop:pet-ch2-endomorphism-adjoint-type-change}
  \lean{PetersenLib.endomorphismDerivation_adjoint_skew}
  \leanok
  \uses{def:pet-ch2-endomorphism-derivation,def:pet-ch2-tensor-algebra-inner-product}
  Assume $V$ has an inner product. (1) The adjoint of $L:V\to V$ extends to
  the adjoint of $L:T(V)\to T(V)$. (2) If $L\in\mathfrak{so}(V)$ (skew-adjoint),
  then $L$ commutes with type change of tensors.
  \source{petersen:page-0082}
\end{proposition}

\subsection{Derivatives}

\begin{proposition}[Prop. 2.3.4 — $L_U$ decomposes through $\nabla$]
  \label{prop:pet-ch2-lie-derivative-as-endomorphism-derivation}
  \lean{PetersenLib.lieDerivative_eq_covariant_minus_endomorphism}
  \leanok
  \uses{def:pet-ch2-lie-derivative-vector-field,def:pet-ch2-covariant-derivative-tensor,def:pet-ch2-endomorphism-derivation}
  Viewing $\nabla U$ as the $(1,1)$-tensor $X\mapsto\nabla_XU$, the Lie
  derivative in the direction $U$ decomposes as $L_U=\nabla_U-(\nabla U)$,
  where $(\nabla U)$ acts by the endomorphism derivation of
  \cref{def:pet-ch2-endomorphism-derivation}.
  \source{petersen:page-0082,petersen:page-0083}
\end{proposition}
\begin{proof}
  \uses{prop:pet-ch2-lie-derivative-as-endomorphism-derivation}
  \leanok
  Both sides are derivations on tensors, so it suffices to check the identity
  on functions and on vector fields. On functions, $L_Uf=D_Uf=\nabla_Uf$ while
  $(\nabla U)$ acts as zero on scalars, matching $\nabla_U-(\nabla U)$. On a
  vector field $Y$, $L_UY=[U,Y]=\nabla_UY-\nabla_YU$
  (\cref{prop:pet-ch2-lie-derivative-bracket}, torsion-freeness), and
  $(\nabla U)(Y)=\nabla_YU$ by definition of the induced endomorphism action
  on vectors (\cref{def:pet-ch2-endomorphism-derivation}), so
  $\big(\nabla_U-(\nabla U)\big)(Y)=\nabla_UY-\nabla_YU=L_UY$.
\end{proof}

\section{The Connection in Tensor Notation}

\begin{definition}[coordinate formulas for the metric]
  \label{def:pet-ch2-metric-coordinate-formulas}
  \lean{PetersenLib.metricCoordinateFormulas}
  \leanok
  In local coordinates $g=g_{ij}dx^idx^j$; for $X=X^i\partial_i$, $Y=Y^j\partial_j$,
  $g(X,Y)=g_{ij}X^iY^j$, and the dual $1$-form (\cref{def:pet-ch2-dual-one-form})
  is $\theta_X=g_{ij}dx^i(X)dx^j=g_{ij}X^idx^j$. The inverse matrix $[g^{ij}]$ of
  $[g_{ij}]$ satisfies $\delta^i_j=g^{ik}g_{kj}$.
  \source{petersen:page-0083}
\end{definition}

\begin{definition}[vector field dual to a $1$-form]
  \label{def:pet-ch2-vector-dual-to-form}
  \lean{PetersenLib.vectorDualToForm}
  \leanok
  \uses{def:pet-ch2-metric-coordinate-formulas}
  The vector field $X$ dual to $\omega=\omega_idx^i$, defined implicitly by
  $g(X,Y)=\omega(Y)$, is $X=g^{ij}\omega_j\partial_i$: from $\theta_X=g_{ij}X^idx^j=\omega_jdx^j=\omega$
  one gets $g_{ij}X^i=\omega_j$, and multiplying by $g^{kj}$ and using symmetry
  of $g_{ij}$ gives $g^{kj}\omega_j=g^{kj}g_{ij}X^i=\delta^k_iX^i=X^k$.
  \source{petersen:page-0083}
\end{definition}

\begin{remark}[gradient in coordinates]
  \label{rem:pet-ch2-gradient-coordinate-formula}
  \lean{PetersenLib.gradient_coordinate_formula}
  \leanok
  \uses{def:pet-ch2-directional-derivative,def:pet-ch2-vector-dual-to-form}
  Applying \cref{def:pet-ch2-vector-dual-to-form} to $\omega=df=\partial_jf\,dx^j$
  gives $\nabla f=g^{ij}\partial_jf\,\partial_i$, matching
  \cref{def:pet-ch2-directional-derivative}.
  \source{petersen:page-0084}
\end{remark}

\begin{definition}[Christoffel symbols of the second kind]
  \label{def:pet-ch2-christoffel-symbols-second-kind}
  \lean{PetersenLib.christoffelSymbolsSecondKind}
  \leanok
  \uses{def:pet-ch2-covariant-derivative-tensor,def:pet-ch2-metric-coordinate-formulas}
  In local coordinates, expand
  $\nabla_YX=Y^i\big(\partial_i(X^j)\big)\partial_j+Y^iX^j\Gamma^k_{ij}\partial_k$, where
  $\Gamma^k_{ij}$, the \emph{Christoffel symbols of the second kind}, are
  defined by $\nabla_{\partial_i}\partial_j=\Gamma^k_{ij}\partial_k$.
  \source{petersen:page-0084}
\end{definition}

\begin{proposition}[$\Gamma^k_{ij}$ in terms of the metric]
  \label{prop:pet-ch2-christoffel-formula}
  \lean{PetersenLib.christoffelSymbols_metric_formula}
  \leanok
  \uses{def:pet-ch2-christoffel-symbols-second-kind,rem:pet-ch2-koszul-formula}
  $\displaystyle\Gamma^k_{ij}=\tfrac12g^{k\ell}\big(\partial_ig_{\ell j}+\partial_jg_{\ell i}-\partial_\ell g_{ij}\big)$.
  \source{petersen:page-0084,petersen:page-0085}
\end{proposition}
\begin{proof}
  \uses{prop:pet-ch2-christoffel-formula}
  \leanok
  Apply Koszul's formula (\cref{rem:pet-ch2-koszul-formula}) to
  $X=\partial_i,Y=\partial_j,Z=\partial_\ell$, using that all Lie brackets of
  coordinate fields vanish: the left-hand side is
  $2g(\nabla_{\partial_j}\partial_i,\partial_\ell)=2g(\Gamma^k_{ij}\partial_k,\partial_\ell)=2\Gamma^k_{ij}g_{k\ell}$,
  and the right-hand side reduces to
  $\partial_ig_{j\ell}+\partial_jg_{i\ell}-\partial_\ell g_{ij}$. Multiplying by
  $g^{\ell m}$ and summing over $\ell$,
  \[
    2\Gamma^m_{ij}=2\Gamma^k_{ij}g_{k\ell}g^{\ell m}=(\partial_ig_{j\ell}+\partial_jg_{i\ell}-\partial_\ell g_{ij})g^{\ell m},
  \]
  giving $\Gamma^m_{ij}=\tfrac12g^{\ell m}(\partial_ig_{j\ell}+\partial_jg_{i\ell}-\partial_\ell g_{ij})$, which is the
  claimed formula after renaming indices.
\end{proof}

\begin{definition}[Christoffel symbols of the first kind]
  \label{def:pet-ch2-christoffel-symbols-first-kind}
  \lean{PetersenLib.christoffelSymbolsFirstKind}
  \leanok
  \uses{def:pet-ch2-christoffel-symbols-second-kind}
  $\Gamma_{ij,k}:=\tfrac12(\partial_jg_{ik}+\partial_ig_{jk}-\partial_kg_{ij})=g(\nabla_{\partial_i}\partial_j,\partial_k)$;
  classically also written $[ij,k]$, with $\Gamma^k_{ij}=\begin{Bmatrix}k\\i,j\end{Bmatrix}$.
  \source{petersen:page-0085}
\end{definition}

\begin{remark}[Christoffel symbols are not tensorial]
  \label{rem:pet-ch2-christoffel-non-tensorial}
  \lean{PetersenLib.christoffelSymbols_not_tensorial}
  \uses{def:pet-ch2-christoffel-symbols-first-kind}
  The Christoffel symbols may vanish in one coordinate system but not
  another, so they are not tensorial. In the plane, they vanish in Cartesian
  coordinates but not in polar coordinates: $\Gamma_{\theta\theta,r}
  =\tfrac12(\partial_\theta g_{\theta r}+\partial_\theta g_{\theta r}-\partial_rg_{\theta\theta})
  =-\tfrac12\partial_r(r^2)=-r$, since $g_{\theta\theta}=r^2$ and $g_{\theta r}=0$.
  \source{petersen:page-0085}
\end{remark}

\begin{remark}[normal coordinates and the Christoffel symbols]
  \label{rem:pet-ch2-christoffel-normal-coordinates}
  \lean{PetersenLib.christoffelSymbols_vanish_normalCoordinates}
  \uses{def:pet-ch2-normal-coordinates-frame,def:pet-ch2-christoffel-symbols-second-kind}
  In coordinates normal at $p$ (\cref{def:pet-ch2-normal-coordinates-frame}),
  $g_{ij}|_p=\delta_{ij}$ and $\Gamma^k_{ij}|_p=0$; consequently, at $p$, the
  covariant derivative is computed exactly as in Euclidean space:
  $(\nabla_YX)|_p=\big(Y^i\partial_i(X^j)\big)|_p\,\partial_j|_p$.
  \source{petersen:page-0085}
\end{remark}

\begin{proposition}[torsion-free and metric properties in coordinates]
  \label{prop:pet-ch2-christoffel-symmetry-metric-property}
  \lean{PetersenLib.christoffel_symmetric_metric_property}
  \leanok
  \uses{def:pet-ch2-christoffel-symbols-first-kind,def:pet-ch2-christoffel-symbols-second-kind}
  Torsion-freeness (\cref{def:pet-ch2-riemannian-connection}(3)) is equivalent
  to $\Gamma^k_{ij}=\Gamma^k_{ji}$, and the metric property
  (\cref{def:pet-ch2-riemannian-connection}(4)) becomes
  $\partial_kg_{ij}=\Gamma_{ki,j}+\Gamma_{kj,i}$.
  \source{petersen:page-0086}
\end{proposition}
\begin{proof}
  \uses{prop:pet-ch2-christoffel-symmetry-metric-property}
  \leanok
  Since $[\partial_i,\partial_j]=0$, torsion-freeness gives
  $\Gamma^k_{ij}\partial_k=\nabla_{\partial_i}\partial_j=\nabla_{\partial_j}\partial_i=\Gamma^k_{ji}\partial_k$,
  i.e.\ $\Gamma^k_{ij}=\Gamma^k_{ji}$. The metric property applied to
  $\partial_kg_{ij}=D_{\partial_k}g(\partial_i,\partial_j)$ gives
  $\partial_kg_{ij}=g(\nabla_{\partial_k}\partial_i,\partial_j)+g(\partial_i,\nabla_{\partial_k}\partial_j)
  =\Gamma_{ki,j}+\Gamma_{kj,i}$ by \cref{def:pet-ch2-christoffel-symbols-first-kind}.
\end{proof}

\begin{proposition}[Hessian in Christoffel-symbol coordinates]
  \label{prop:pet-ch2-hessian-coordinate-formula}
  \lean{PetersenLib.hessian_coordinate_formula}
  \uses{def:pet-ch2-hessian-lie-derivative,def:pet-ch2-christoffel-symbols-second-kind}
  $\operatorname{Hess}f(\partial_i,\partial_j)=\partial_i\partial_jf-\Gamma^k_{ij}\partial_kf$.
  \source{petersen:page-0086,petersen:page-0087}
\end{proposition}
\begin{proof}
  \uses{prop:pet-ch2-hessian-coordinate-formula}
  By \cref{def:pet-ch2-hessian-lie-derivative} and
  \cref{prop:pet-ch2-lie-derivative-tensor-formula} applied to
  $g=g_{ij}dx^idx^j$ with $X=\nabla f=g^{k\ell}(\partial_kf)\partial_\ell$,
  \[
    2\operatorname{Hess}f(\partial_i,\partial_j)
    = D_{\nabla f}g_{ij}-g(L_{\nabla f}\partial_i,\partial_j)-g(\partial_i,L_{\nabla f}\partial_j).
  \]
  By \cref{prop:pet-ch2-lie-derivative-scaling} applied to the bracket
  $[\nabla f,\partial_i]$, $L_{\nabla f}\partial_i=-\partial_i(g^{k\ell}\partial_kf)\partial_\ell$,
  and $D_{\nabla f}g_{ij}=g^{k\ell}(\partial_kf)(\partial_\ell g_{ij})$. Substituting and
  using the metric-inverse identity
  $0=\partial_k\delta^i_j=(\partial_kg^{i\ell})g_{\ell j}+g^{i\ell}(\partial_kg_{\ell j})$
  to rewrite $\partial_k(g^{ij})$ in terms of $\partial_kg_{ij}$, the expression
  collapses (after collecting the resulting terms and applying
  \cref{prop:pet-ch2-christoffel-formula}) to
  \[
    2\operatorname{Hess}f(\partial_i,\partial_j)
    = 2\partial_i\partial_jf-g^{k\ell}\big(\partial_ig_{\ell j}+\partial_jg_{\ell i}-\partial_\ell g_{ij}\big)\partial_kf
    = 2\big(\partial_i\partial_jf-\Gamma^k_{ij}\partial_kf\big),
  \]
  matching \cref{prop:pet-ch2-christoffel-formula}.
\end{proof}

\begin{definition}[index notation for the covariant derivative]
  \label{def:pet-ch2-covariant-derivative-index-notation}
  \lean{PetersenLib.covariantDerivativeIndexNotation}
  \uses{def:pet-ch2-covariant-derivative-tensor}
  Writing a $(1,k)$-tensor $S=S^i_{j_1\cdots j_k}E_i\otimes\sigma^{j_1}\otimes\cdots\otimes\sigma^{j_k}$
  in a frame $E_i$ with coframe $\sigma^j$, its covariant derivative
  $\nabla S=S^i_{j_1\cdots j_kj_{k+1}}E_i\otimes\sigma^{j_1}\otimes\cdots\otimes\sigma^{j_{k+1}}$
  has coefficients computed by expanding
  $\nabla_{E_{k+1}}S=D_{E_{k+1}}(S^i_{j_1\cdots j_k})E_i\otimes\sigma^{j_1}\otimes\cdots\otimes\sigma^{j_k}
  +S^i_{j_1\cdots j_k}\nabla_{E_{k+1}}(E_i\otimes\sigma^{j_1}\otimes\cdots\otimes\sigma^{j_k})$
  and the product rule on each factor of the frame/coframe tensor product. The
  index-first convention $\nabla_jS:=\nabla_{E_j}S$,
  $\nabla_jS^i_{j_1\cdots j_k}:=(\nabla S)^i_{j_1\cdots j_kj}$, avoids placing
  the differentiation variable last.
  \source{petersen:page-0087,petersen:page-0088}
\end{definition}

\section{Exercises}

\begin{remark}[Exercise 2.5.1 — uniqueness of the Euclidean connection]
  \label{rem:pet-ch2-ex-1}
  \lean{PetersenLib.exercise2_5_1}
  \leanok
  \uses{def:pet-ch2-covariant-derivative-euclidean,def:pet-ch2-affine-connection}
  Show that the connection on Euclidean space is the only affine connection
  such that $\nabla X=0$ for all constant vector fields $X$.
  \source{petersen:page-0088}
\end{remark}

\begin{remark}[Exercise 2.5.2 — regularity needed for skew-symmetry and Jacobi]
  \label{rem:pet-ch2-ex-2}
  \lean{PetersenLib.exercise2_5_2}
  Show that the skew-symmetry property $[X,Y]=-[Y,X]$ does not necessarily
  hold for $C^1$ vector fields, and that the Jacobi identity holds for $C^2$
  vector fields.
  \source{petersen:page-0088}
\end{remark}

\begin{remark}[Exercise 2.5.3 — the torsion tensor]
  \label{rem:pet-ch2-ex-3}
  \lean{PetersenLib.exercise2_5_3}
  \leanok
  \uses{def:pet-ch2-affine-connection}
  Let $\nabla$ be an affine connection. Show that
  $T(X,Y)=\nabla_XY-\nabla_YX-[X,Y]$ defines a $(2,1)$-tensor.
  \source{petersen:page-0088}
\end{remark}

\begin{remark}[Exercise 2.5.4 — extending a tangent to a field]
  \label{rem:pet-ch2-ex-4}
  \lean{PetersenLib.exercise2_5_4}
  Show that if $c:I\to M$ has nonzero speed at $t_0\in I$, there is a vector
  field $X$ with $X|_{c(t)}=\dot c(t)$ for $t$ near $t_0$.
  \source{petersen:page-0088}
\end{remark}

\begin{remark}[Exercise 2.5.5 — product rules for divergence, Laplacian, Hessian]
  \label{rem:pet-ch2-ex-5}
  \lean{PetersenLib.exercise2_5_5}
  \uses{def:pet-ch2-divergence-lie-derivative,def:pet-ch2-laplacian,def:pet-ch2-hessian-lie-derivative}
  For $f,h$ functions and $X$ a vector field on $(M,g)$, show
  $\operatorname{div}(fX)=D_Xf+f\operatorname{div}X$,
  $\Delta(fh)=h\Delta f+f\Delta h+2g(\nabla f,\nabla h)$, and
  $\operatorname{Hess}(fh)=h\operatorname{Hess}f+f\operatorname{Hess}h+df\otimes dh+dh\otimes df$.
  \source{petersen:page-0088}
\end{remark}

\begin{remark}[Exercise 2.5.6 — chain rule for Laplacian and Hessian]
  \label{rem:pet-ch2-ex-6}
  \lean{PetersenLib.exercise2_5_6}
  \uses{def:pet-ch2-laplacian,def:pet-ch2-hessian-lie-derivative}
  For $f$ on $(M,g)$ and $\phi:\mathbb{R}\to\mathbb{R}$, show
  $\Delta(\phi(f))=\dot\phi(f)\Delta f+\ddot\phi(f)|df|^2$ and
  $\operatorname{Hess}(\phi(f))=\dot\phi(f)\operatorname{Hess}f+\ddot\phi(f)\,df^2$.
  \source{petersen:page-0089}
\end{remark}

\begin{remark}[Exercise 2.5.7 — $d\theta_X$ via the connection]
  \label{rem:pet-ch2-ex-7}
  \lean{PetersenLib.exercise2_5_7}
  \leanok
  \uses{def:pet-ch2-dual-one-form,def:pet-ch2-covariant-derivative-tensor}
  For a vector field $X$ with dual $1$-form $\theta_X$, show
  $d\theta_X(Y,Z)=g(\nabla_YX,Z)-g(Y,\nabla_ZX)$.
  \source{petersen:page-0089}
\end{remark}

\begin{remark}[Exercise 2.5.8 — derivatives of the inverse metric coefficients]
  \label{rem:pet-ch2-ex-8}
  \lean{PetersenLib.exercise2_5_8}
  \uses{def:pet-ch2-metric-coordinate-formulas,def:pet-ch2-christoffel-symbols-second-kind}
  Show that $\partial_kg^{ij}$ can be written in terms of $g^{i\ell}$, $g^{j\ell}$
  and the Christoffel symbols; specifically
  $\partial_kg^{ij}=-g^{i\ell}\Gamma^j_{k\ell}-g^{j\ell}\Gamma^i_{k\ell}$.
  \source{petersen:page-0089}
\end{remark}

\begin{remark}[Exercise 2.5.9 — covariant differentiation commutes with contraction and type change]
  \label{rem:pet-ch2-ex-9}
  \lean{PetersenLib.exercise2_5_9}
  \uses{def:pet-ch2-covariant-derivative-tensor}
  For a vector field $X$: (1) show $\operatorname{tr}(\nabla_XS)=\nabla_X(\operatorname{tr}S)$
  for a $(1,1)$-tensor $S$; (2) for $T(Y,Z)=g(S(Y),Z)$, show
  $(\nabla_XT)(Y,Z)=g((\nabla_XS)(Y),Z)$; (3) show contraction and covariant
  differentiation commute in general; (4) show type change and covariant
  differentiation commute.
  \source{petersen:page-0089}
\end{remark}

\begin{remark}[Exercise 2.5.10 — Lie differentiation commutes with contraction]
  \label{rem:pet-ch2-ex-10}
  \lean{PetersenLib.exercise2_5_10}
  \uses{def:pet-ch2-lie-derivative-tensor}
  For a vector field $X$: (1) show $\operatorname{tr}(L_XS)=L_X(\operatorname{tr}S)$
  for a $(1,1)$-tensor $S$; (2) for $T(Y,Z)=g(S(Y),Z)$, show
  $(L_XT)(Y,Z)=(L_Xg)(S(Y),Z)+g((L_XS)(Y),Z)$; (3) show contraction and Lie
  differentiation commute.
  \source{petersen:page-0089}
\end{remark}

\begin{remark}[Exercise 2.5.11 — local gradient fields via self-adjointness]
  \label{rem:pet-ch2-ex-11}
  \lean{PetersenLib.exercise2_5_11}
  \uses{def:pet-ch2-covariant-derivative-tensor}
  Show that $X$ is locally a gradient field iff $Z\mapsto\nabla_ZX$ is
  self-adjoint.
  \source{petersen:page-0089}
\end{remark}

\begin{remark}[Exercise 2.5.12 — isometries and pushforward of the connection]
  \label{rem:pet-ch2-ex-12}
  \lean{PetersenLib.exercise2_5_12}
  \uses{def:pet-ch2-covariant-derivative-tensor}
  For $F:M\to M$ a diffeomorphism, the pushforward of a vector field is
  $(F_*X)|_p=DF(X|_{F^{-1}(p)})$. Let $F$ be an isometry on $(M,g)$: (1) show
  $F_*(\nabla_YX)=\nabla_{F_*Y}F_*X$ for all vector fields; (2) deduce that
  isometries of $(\mathbb{R}^n,g_{\mathrm{eu}})$ are $F(x)=Ox+b$ with $O\in O(n)$,
  $b\in\mathbb{R}^n$, by showing $F$ maps constant vector fields to constant vector
  fields.
  \source{petersen:page-0089,petersen:page-0090}
\end{remark}

\begin{remark}[Exercise 2.5.13 — affine vector fields]
  \label{rem:pet-ch2-ex-13}
  \lean{PetersenLib.exercise2_5_13}
  \uses{def:pet-ch2-lie-derivative-connection,def:pet-ch2-killing-field}
  A vector field $X$ is \emph{affine} if $L_X\nabla=0$. (1) Show Killing
  fields are affine; (2) give an example of an affine field on $\mathbb{R}^n$ that
  is not Killing.
  \source{petersen:page-0090}
\end{remark}

\begin{remark}[Exercise 2.5.14 — the left-invariant connection on a Lie group]
  \label{rem:pet-ch2-ex-14}
  \lean{PetersenLib.exercise2_5_14}
  \uses{def:pet-ch2-affine-connection}
  For a Lie group $G$, show there is a unique affine connection with
  $\nabla X=0$ for all left-invariant vector fields $X$, and that this
  connection is torsion free iff the Lie algebra is Abelian.
  \source{petersen:page-0090}
\end{remark}

\begin{remark}[Exercise 2.5.15 — chain rule for the Hessian]
  \label{rem:pet-ch2-ex-15}
  \lean{PetersenLib.exercise2_5_15}
  \uses{def:pet-ch2-hessian-lie-derivative}
  Show $\operatorname{Hess}(\phi(f))=\phi''\,df^2+\phi'\operatorname{Hess}f$.
  \source{petersen:page-0090}
\end{remark}

\begin{remark}[Exercise 2.5.16 — derivations are $L_X+L$]
  \label{rem:pet-ch2-ex-16}
  \lean{PetersenLib.exercise2_5_16}
  \uses{def:pet-ch2-lie-derivative-tensor,def:pet-ch2-tensor-derivation,def:pet-ch2-endomorphism-derivation}
  For a vector field $X$ and $(1,1)$-tensor $L$: (1) show $L_X+L$ is a
  derivation on tensors; (2) show every derivation has this form with $X$
  unique; (3) show derivations are determined by their action on functions and
  vector fields; (4) show $L_{fX}=fL_X-X\otimes df$, where $X\otimes df$ is
  the rank-$1$ $(1,1)$-tensor $Y\mapsto X\,df(Y)$.
  \source{petersen:page-0090}
\end{remark}

\begin{remark}[Exercise 2.5.17 — constant-length fields are orthogonal to their covariant derivative]
  \label{rem:pet-ch2-ex-17}
  \lean{PetersenLib.exercise2_5_17}
  \leanok
  \uses{def:pet-ch2-covariant-derivative-tensor}
  Show that if $X$ has constant length on $(M,g)$, then $\nabla_ZX\perp X$ for
  all $Z$.
  \source{petersen:page-0090}
\end{remark}

\begin{remark}[Exercise 2.5.18 — Lie and covariant derivatives agree where $T$ vanishes]
  \label{rem:pet-ch2-ex-18}
  \lean{PetersenLib.exercise2_5_18}
  \uses{def:pet-ch2-lie-derivative-tensor,def:pet-ch2-covariant-derivative-tensor,def:pet-ch2-hessian-lie-derivative}
  Show that if a tensor field $T$ on $(M,g)$ vanishes at $p$, then
  $L_XT|_p=\nabla_XT|_p$ for any vector field $X$; conclude that the $(1,1)$
  version of the Hessian of a function is independent of the metric at a
  critical point, and interpret $L_XT|_p$.
  \source{petersen:page-0090}
\end{remark}

\begin{remark}[Exercise 2.5.19 — existence of a normal frame]
  \label{rem:pet-ch2-ex-19}
  \lean{PetersenLib.exercise2_5_19}
  \uses{def:pet-ch2-normal-coordinates-frame}
  For $p\in(M,g)$ and orthonormal $e_1,\dots,e_n\in T_pM$, show there is an
  orthonormal frame $E_1,\dots,E_n$ near $p$ with $E_i(p)=e_i$ and
  $(\nabla E_i)|_p=0$.
  \source{petersen:page-0090}
\end{remark}

\begin{remark}[Exercise 2.5.20 — existence of normal coordinates]
  \label{rem:pet-ch2-ex-20}
  \lean{PetersenLib.exercise2_5_20}
  \uses{def:pet-ch2-normal-coordinates-frame}
  Show there are coordinates $x^1,\dots,x^n$ at $p$ with $\partial_i=e_i$ and
  $(\nabla\partial_i)|_p=0$, so $g_{ij}|_p=\delta_{ij}$ and
  $\partial_kg_{ij}|_p=0$ — \emph{normal coordinates} at $p$.
  \source{petersen:page-0090,petersen:page-0091}
\end{remark}

\begin{remark}[Exercise 2.5.21 — transformation of Christoffel symbols]
  \label{rem:pet-ch2-ex-21}
  \lean{PetersenLib.exercise2_5_21}
  \uses{def:pet-ch2-christoffel-symbols-second-kind,def:pet-ch2-christoffel-symbols-first-kind}
  For coordinates $x^i,\tilde x^i$ at $p$, show the Christoffel symbols
  transform by
  \[
    \tilde\Gamma^k_{ij}
    = \frac{\partial^2x^s}{\partial\tilde x^i\partial\tilde x^j}\frac{\partial\tilde x^k}{\partial x^s}
    + \frac{\partial x^s}{\partial\tilde x^i}\frac{\partial x^t}{\partial\tilde x^j}\frac{\partial\tilde x^k}{\partial x^\ell}\Gamma^\ell_{st},
  \]
  equivalently, isolating the second-partial term,
  $\dfrac{\partial^2x^r}{\partial\tilde x^i\partial\tilde x^j}
  = \Gamma^k_{ij}\dfrac{\partial x^r}{\partial\tilde x^k}
  -\dfrac{\partial x^r}{\partial\tilde x^k}\dfrac{\partial x^s}{\partial\tilde x^\ell}\tilde\Gamma^k_{ij}$
  in the appropriately contracted sense, and
  \[
    \tilde\Gamma_{ij,k}
    = \frac{\partial^2x^s}{\partial\tilde x^i\partial\tilde x^j}\frac{\partial x^\ell}{\partial\tilde x^k}g_{s\ell}
    + \frac{\partial x^s}{\partial\tilde x^i}\frac{\partial x^t}{\partial\tilde x^j}\frac{\partial x^\ell}{\partial\tilde x^k}\Gamma_{st,\ell}.
  \]
  \source{petersen:page-0091}
\end{remark}

\begin{remark}[Exercise 2.5.22 — Christoffel symbols of an induced metric]
  \label{rem:pet-ch2-ex-22}
  \lean{PetersenLib.exercise2_5_22}
  \uses{def:pet-ch2-metric-coordinate-formulas,def:pet-ch2-christoffel-symbols-first-kind}
  For $M^n\subset\mathbb{R}^{n+m}$ with induced metric and parametrization
  $u^s(x^1,\dots,x^n)$, $s=1,\dots,n+m$, show
  $g_{ij}=\sum_su^s_iu^s_j$ and
  $\Gamma_{ij,k}=\sum_su^s_ku^s_{ij}$, where $u^s_i=\partial u^s/\partial x^i$
  and $u^s_{ij}=\partial^2u^s/\partial x^i\partial x^j$.
  \source{petersen:page-0091}
\end{remark}

\begin{remark}[Exercise 2.5.23 — volume form and Laplacian in coordinates]
  \label{rem:pet-ch2-ex-23}
  \lean{PetersenLib.exercise2_5_23}
  \uses{def:pet-ch2-metric-coordinate-formulas,def:pet-ch2-laplacian}
  On an oriented $(M,g)$: (1) if $v_1,\dots,v_n$ is positively oriented, show
  $\mathrm{vol}(v_1,\dots,v_n)=\sqrt{\det(g(v_i,v_j))}$; (2) in positively
  oriented coordinates, show
  $\mathrm{vol}=\sqrt{\det(g_{ij})}\,dx^1\wedge\cdots\wedge dx^n$; (3) conclude
  $\Delta u=\frac1{\sqrt{\det(g_{ij})}}\partial_k\big(\sqrt{\det(g_{ij})}\,g^{k\ell}\partial_\ell u\big)$,
  which at a point with normal coordinates reduces to $\Delta f(p)=\sum_i\partial_i^2f$.
  \source{petersen:page-0091,petersen:page-0092}
\end{remark}

\begin{remark}[Exercise 2.5.24 — covariant derivative components of $(0,2)$- and $(1,1)$-tensors]
  \label{rem:pet-ch2-ex-24}
  \lean{PetersenLib.exercise2_5_24}
  \uses{def:pet-ch2-christoffel-symbols-second-kind,def:pet-ch2-covariant-derivative-index-notation}
  Show that a $(0,2)$-tensor $T_{ik}$ has
  $(\nabla T)_{jk\ell}=\partial_jT_{k\ell}-\Gamma^m_{jk}T_{m\ell}-\Gamma^m_{j\ell}T_{km}$,
  and a $(1,1)$-tensor $T^k_i$ has
  $(\nabla T)^k_{ji}=\partial_jT^k_i-\Gamma^\ell_{ji}T^k_\ell+\Gamma^k_{j\ell}T^\ell_i$.
  \source{petersen:page-0092}
\end{remark}

\begin{remark}[Exercise 2.5.25 — second fundamental form for isometric immersions]
  \label{rem:pet-ch2-ex-25}
  \lean{PetersenLib.exercise2_5_25}
  \uses{def:pet-ch2-covariant-derivative-tensor}
  Let $F:(M,g_M)\hookrightarrow(\tilde M,g_{\tilde M})$ be an isometric
  immersion. Show that the tangential component of
  $\nabla^{\tilde M}_XY$ (for $X,Y$ tangent to $M$) is $\nabla^M_XY$, that the
  normal component $\mathrm{I\!I}(X,Y)=\nabla^{\tilde M}_XY-\nabla^M_XY$ is
  symmetric in $X,Y$, and conclude $\mathrm{I\!I}$ is tensorial.
  \source{petersen:page-0092}
\end{remark}

\begin{remark}[Exercise 2.5.26 — normal connection along an immersed submanifold]
  \label{rem:pet-ch2-ex-26}
  \lean{PetersenLib.exercise2_5_26}
  \uses{def:pet-ch2-covariant-derivative-tensor}
  With $F$ as in \cref{rem:pet-ch2-ex-25} and $T^\perp M$ the normal bundle,
  for a normal field $V$ along $M$ and tangent field $X$: (1) show
  $\nabla^{\tilde M}_XV$ is defined; (2) decompose it as
  $\nabla^{\tilde M}_XV=\nabla^\perp_XV+T_XV$ into normal and tangential
  parts, and show $g_M(T_XY,V)=-g_M(Y,T_XV)$; (3) show $\nabla^\perp_XV$ is
  linear and a derivation in $V$, tensorial in $X$.
  \source{petersen:page-0092,petersen:page-0093}
\end{remark}

\begin{remark}[Exercise 2.5.27 — divergence theorem and the maximum principle]
  \label{rem:pet-ch2-ex-27}
  \lean{PetersenLib.exercise2_5_27}
  \uses{def:pet-ch2-laplacian,def:pet-ch2-divergence-lie-derivative}
  On an oriented $(M,g)$: (1) if $f$ has compact support,
  $\int_M\Delta f\cdot\mathrm{vol}=0$; (2)
  $\operatorname{div}(fX)=g(\nabla f,X)+f\operatorname{div}X$; (3)
  $\Delta(f_1f_2)=(\Delta f_1)f_2+2g(\nabla f_1,\nabla f_2)+f_1\Delta f_2$; (4)
  Green's formula
  $\int_Mf_1\Delta f_2\,\mathrm{vol}=-\int_Mg(\nabla f_1,\nabla f_2)\,\mathrm{vol}$
  for compactly supported $f_1$; (5) conclude the weak maximum principle: a
  compactly supported subharmonic or superharmonic function is constant; more
  generally a subharmonic (superharmonic) function attaining a global maximum
  (minimum) is constant, without a compact-support hypothesis (the strong
  maximum principle).
  \source{petersen:page-0093}
\end{remark}

\begin{remark}[Exercise 2.5.28 — incompressible vector fields]
  \label{rem:pet-ch2-ex-28}
  \lean{PetersenLib.exercise2_5_28}
  \uses{def:pet-ch2-divergence-lie-derivative,def:pet-ch2-covariant-derivative-tensor}
  A vector field $X$ (and its flow) is \emph{incompressible} if
  $\operatorname{div}X=0$. (1) Show $X$ is incompressible iff its local flows
  are volume preserving; (2) show a unit vector field $X$ on $\mathbb{R}^2$ that is
  incompressible has $\nabla X=0$; (3) find an incompressible unit vector
  field $X$ on $\mathbb{R}^3$ with $\nabla X\ne0$.
  \source{petersen:page-0093}
\end{remark}

\begin{remark}[Exercise 2.5.29 — geodesic unit fields and integrability]
  \label{rem:pet-ch2-ex-29}
  \lean{PetersenLib.exercise2_5_29}
  \uses{def:pet-ch2-dual-one-form,def:pet-ch2-covariant-derivative-tensor}
  Let $X$ be a unit vector field on $(M,g)$ with $\nabla_XX=0$. (1) Show $X$
  is locally a gradient field iff the orthogonal distribution $X^\perp$ is
  integrable; (2) show $X^\perp$ is integrable near $p$ if it has an integral
  submanifold through $p$ (hint: show $L_X\theta_X=0$); (3) find such an $X$
  on $S^3$ that is not a gradient field.
  \source{petersen:page-0093}
\end{remark}

\begin{remark}[Exercise 2.5.30 — parallel complementary distributions]
  \label{rem:pet-ch2-ex-30}
  \lean{PetersenLib.exercise2_5_30}
  \uses{def:pet-ch2-covariant-derivative-tensor}
  Let $E,F$ be orthogonal-complement distributions on $(M,g)$, each parallel
  ($X,Y$ tangent to $E$ implies $\nabla_XY$ tangent to $E$, similarly for
  $F$). (1) Show $E,F$ are integrable; (2) show every $p\in M$ has a product
  neighborhood $U=V_E\times V_F$ with
  $(U,g)=(V_E\times V_F,g|_{V_E}+g|_{V_F})$, the integral submanifolds through
  $p$.
  \source{petersen:page-0094}
\end{remark}

\begin{remark}[Exercise 2.5.31 — parallel vector fields and product metrics]
  \label{rem:pet-ch2-ex-31}
  \lean{PetersenLib.exercise2_5_31}
  \uses{def:pet-ch2-parallel-tensor,rem:pet-ch2-ex-30}
  Let $X$ be a parallel vector field on $(M,g)$. Show $X$ has constant
  length, that $X$ generates parallel distributions (one containing $X$, the
  other its orthogonal complement), and conclude that locally
  $(U,g)=(V\times I,g|_V+dt^2)$ with $V$ perpendicular to $X$.
  \source{petersen:page-0094}
\end{remark}

\begin{remark}[Exercise 2.5.32 — Leibniz rule for the inner product of tensors]
  \label{rem:pet-ch2-ex-32}
  \lean{PetersenLib.exercise2_5_32}
  \uses{def:pet-ch2-covariant-derivative-tensor,def:pet-ch2-tensor-algebra-inner-product}
  For tensors $S,T$ of the same type, show
  $D_Xg(S,T)=g(\nabla_XS,T)+g(S,\nabla_XT)$.
  \source{petersen:page-0094}
\end{remark}

\begin{remark}[Exercise 2.5.33 — almost complex structures, Nijenhuis tensor, K\"ahler condition]
  \label{rem:pet-ch2-ex-33}
  \lean{PetersenLib.exercise2_5_33}
  \uses{def:pet-ch2-tensor-derivation,def:pet-ch2-covariant-derivative-tensor,def:pet-ch2-lie-derivative-tensor}
  An \emph{almost complex structure} is a $(1,1)$-tensor $J$ with $J^2=-I$; a
  \emph{Hermitian structure} on $(M,g)$ is such $J$ with
  $g(J(X),J(Y))=g(X,Y)$, and its \emph{K\"ahler form} is
  $\omega(X,Y)=g(J(X),Y)$. (1) Show the \emph{Nijenhuis tensor}
  $N(X,Y)=[J(X),J(Y)]-J([J(X),Y])-J([X,J(Y)])-[X,Y]$ is a tensor; (2) show
  $N=0$ if $J$ comes from a complex structure (the converse is the
  Newlander–Nirenberg theorem); (3) show $\omega$ is a $2$-form; (4) show
  $d\omega=0$ if $\nabla J=0$; (5) conversely show that if $d\omega=0$ and $J$
  is a complex structure, then $\nabla J=0$ — such $(M,g,J)$ is called
  \emph{K\"ahler}.
  \source{petersen:page-0094}
\end{remark}

\begin{remark}[Exercise 2.5.34 — index notation for the covariant derivative]
  \label{rem:pet-ch2-ex-34}
  \lean{PetersenLib.exercise2_5_34}
  \uses{def:pet-ch2-covariant-derivative-index-notation,def:pet-ch2-divergence-lie-derivative,def:pet-ch2-covariant-derivative-adjoint}
  Write $\nabla_iT$ for the covariant derivative in the direction of the
  $i$th coordinate field, and $\nabla^iT=g^{ij}\nabla_jT$. (1) For a function
  $f$, show $df=\nabla_if\,dx^i$ and $\nabla f=\nabla^if\,\partial_i$; (2) for a
  vector field $X$, show $(\nabla_iX)^i=\operatorname{div}X$; (3) for a
  $(0,2)$-tensor $T$, show $(\nabla^iT)_{ij}=-(\nabla^*T)_j$.
  \source{petersen:page-0094}
\end{remark}
